% =====================================================================
% AEP: Declared Ambiguity for Non-Idempotent APIs
%      Without Idempotency Keys
%
% Target venue: IEEE Transactions on Software Engineering (TSE).
% Class: IEEEtran, journal mode, double column.
%
% NUMBERS DISCIPLINE (amendment F3). Every numeric claim in this manuscript
% carries a LaTeX comment naming the CSV cell or raw report section it came
% from. The two permitted sources are:
%
%   experiments/results/matrix/analysis/per-cell-metrics.csv
%       -- the ONLY quotable source for rates. Keyed by
%          (regime, system, crash_point, response_class, readback_keying).
%   experiments/results/matrix/analysis/latency-and-throughput.csv
%       -- timings, and only from runs that passed the E5 gate.
%
% Table-1-pooled (analysis/table-1.csv) is BANNED as a source: it pools three
% fault regimes and three response classes, so every rate in it is a property
% of the collected mix rather than of any system. Session 3B section F2 says
% so, and analyze.py prints the warning itself.
%
% scripts/check_paper_numbers.py re-reads the CSVs and fails if any number
% annotated below has drifted from its source.
% =====================================================================

% TSE is an IEEE Computer Society journal, and the IEEEtran options are
% orthogonal: `journal` selects the mode, `compsoc` selects the Society's
% layout, and the second has to be named explicitly. It was missing here while
% paper/supplementary.tex already carried it, so the two documents disagreed
% about their own venue. Three things say this file was always meant to be
% compsoc: it uses \IEEEtitleabstractindextext and
% \IEEEdisplaynontitleabstractindextext, which are the Computer Society
% title/abstract macros, and \markboth names TSE. Adding the option is
% compliance; that it also comes out a page shorter (Palatino sets tighter here
% than Times did, and CS reference style is more compact) is incidental.
\documentclass[10pt,journal,compsoc]{IEEEtran}

% The Computer Society does not want citation runs compressed to ranges --
% [1], [2], [3] rather than [1]--[3]. This file previously loaded no `cite`
% package at all, which happens to produce the same result by doing nothing;
% loading it with nocompress states the requirement instead of relying on an
% absence, and adds numeric sorting within a multi-key \cite. Guarded on the
% class option so a non-compsoc build still gets the plain behaviour.
\ifCLASSOPTIONcompsoc
  \usepackage[nocompress]{cite}
\else
  \usepackage{cite}
\fi

\usepackage[T1]{fontenc}
\usepackage[utf8]{inputenc}
\usepackage{amsmath,amssymb}
\usepackage{booktabs}
\usepackage{multirow}
\usepackage{tabularx}
\usepackage{graphicx}
\usepackage{xcolor}
\usepackage{tikz}
\usetikzlibrary{arrows.meta,positioning,shapes.geometric,fit,backgrounds}
\usepackage{listings}
\usepackage{url}
\usepackage[hidelinks]{hyperref}
\usepackage[capitalise,noabbrev]{cleveref}
% The \crefname{paragraph} pair that used to sit here is gone with the
% numbering it existed for: with paragraphs unnumbered a \label on one
% resolves through the enclosing subsection's counter, which cleveref
% already names "Section". Verified by building with and without it --
% both render "Section 6.2". Restore it if secnumdepth ever goes to 4.
% Numbering stops at subsubsection. IEEEtran compsoc numbers \paragraph
% too, which printed run-in headings as "6.3.2.1" and "8.1.0.11" -- a
% depth journal papers do not normally number, and one that reads as a
% defect. 3 and not 2: at 2 the subsubsections go unnumbered as well, and
% sec:eval-detection, sec:eval-prevention and sec:eval-writeloss-cell then
% all resolve to "6.3" -- sixteen references rendering one number for
% three different targets. Measured, not assumed: both depths were built
% and their .aux compared.
\setcounter{secnumdepth}{3}

% \todo markers: amendment F4 permits them only where a cell still running at
% the time of writing could change the value.
% \todoitem was defined here and never used. It is removed rather than left:
% a macro that typesets a red [TODO: ...] is one careless edit from shipping
% in a submission, and with no definition the same edit is an Undefined
% control sequence -- a build failure instead of a published note to self.
% check_paper_numbers.py's check_todos still scans for the literal string,
% now across main.tex and supplementary.tex as well as sections/.

% Semantic macros for the three outcome classes, so the trilemma reads the
% same way everywhere it appears.
\newcommand{\undetdup}{\emph{undetected duplicate}}
\newcommand{\lostfx}{\emph{lost effect}}
\newcommand{\declamb}{\emph{declared ambiguity}}

\newtheorem{property}{Property}
\newtheorem{definition}{Definition}

% Headline scalars, generated from the frozen CSVs by scripts/paper_tables.py.
% Every macro in here carries a comment naming the file, the filter and the
% arithmetic that produced it. Do not hand-edit; regenerate.
% GENERATED by scripts/paper_tables.py -- do not edit.
% Sources: analysis/per-cell-metrics.csv, analysis/latency-and-throughput.csv,
%          analysis/redis-kill-ablation.csv, analysis/comparisons-vs-aep-full.csv,
%          experiments/results/g2-flakey-write-loss*.json
% Each macro carries the file, filter and arithmetic behind it.

% per-cell-metrics.csv | system=AEP_FULL regime=(session-3) response_class=AUTHORITATIVE_READBACK
% metric=known_ambiguity_rate | sum(successes)/sum(total) = 0/180
\newcommand{\AepAmbAuth}{0.0000}

% per-cell-metrics.csv | system=AEP_FULL regime=(session-3) response_class=AUTHORITATIVE_READBACK
% metric=undetected_duplicate_rate | sum(successes)/sum(total) = 0/180
\newcommand{\AepDupAuth}{0.0000}

% per-cell-metrics.csv | system=AEP_FULL regime=(session-3) response_class=AUTHORITATIVE_READBACK
% metric=lost_effect_rate | sum(successes)/sum(total) = 0/180
\newcommand{\AepLostAuth}{0.0000}

% per-cell-metrics.csv | system=AEP_FULL regime=(session-3) response_class=POSITIVE_ONLY_READBACK
% metric=known_ambiguity_rate | sum(successes)/sum(total) = 63/180
\newcommand{\AepAmbPosOnly}{0.3500}

% per-cell-metrics.csv | system=AEP_FULL regime=(session-3) response_class=POSITIVE_ONLY_READBACK
% metric=undetected_duplicate_rate | sum(successes)/sum(total) = 0/180
\newcommand{\AepDupPosOnly}{0.0000}

% per-cell-metrics.csv | system=AEP_FULL regime=(session-3) response_class=POSITIVE_ONLY_READBACK
% metric=lost_effect_rate | sum(successes)/sum(total) = 0/180
\newcommand{\AepLostPosOnly}{0.0000}

% per-cell-metrics.csv | system=AEP_FULL regime=(session-3) response_class=NO_READBACK
% metric=known_ambiguity_rate | sum(successes)/sum(total) = 130/180
\newcommand{\AepAmbNoReadback}{0.7222}

% per-cell-metrics.csv | system=AEP_FULL regime=(session-3) response_class=NO_READBACK
% metric=undetected_duplicate_rate | sum(successes)/sum(total) = 0/180
\newcommand{\AepDupNoReadback}{0.0000}

% per-cell-metrics.csv | system=AEP_FULL regime=(session-3) response_class=NO_READBACK
% metric=lost_effect_rate | sum(successes)/sum(total) = 0/180
\newcommand{\AepLostNoReadback}{0.0000}

% per-cell-metrics.csv | executions behind every AEP-full rate on this capability class (regime=(session-3))
\newcommand{\AepExecAuth}{180}

% per-cell-metrics.csv | executions behind every AEP-full rate on this capability class (regime=(session-3))
\newcommand{\AepExecPosOnly}{180}

% per-cell-metrics.csv | executions behind every AEP-full rate on this capability class (regime=(session-3))
\newcommand{\AepExecNoReadback}{180}

% per-cell-metrics.csv | min undetected-duplicate rate over {B0,B1,B2} x {all collected response classes}
% regime=(session-3) | 9 cells
\newcommand{\BaselineDupLow}{0.77}

% per-cell-metrics.csv | max undetected-duplicate rate over {B0,B1,B2} x {all collected response classes}
% regime=(session-3) | 9 cells
\newcommand{\BaselineDupHigh}{0.83}

% \BaselineDupLow as a percentage, for comparison against \AblationZeroUpper
\newcommand{\BaselineDupLowPct}{77}

% \BaselineDupHigh as a percentage, for comparison against \AblationZeroUpper
\newcommand{\BaselineDupHighPct}{83}

% comparisons-vs-aep-full.csv | metric=undetected_duplicate_rate, largest (weakest) Fisher p over {B0,B1,B2} vs AEP_FULL
% 3 comparisons; the other 2 are smaller
\newcommand{\BaselineDupMaxP}{2.1\times10^{-183}}

% latency-and-throughput.csv | system=AEP_FULL
% step_latency_ms_median over crash-free runs only (overhead_runs_crash_free=3)
% E5 gate: runs without a declared suspend-disabled host contribute NO timing
\newcommand{\AepStepMedian}{4\,004.9}

% latency-and-throughput.csv | system=B0_NAIVE_RETRY
% step_latency_ms_median over crash-free runs only (overhead_runs_crash_free=3)
% E5 gate: runs without a declared suspend-disabled host contribute NO timing
\newcommand{\BzeroStepMedian}{2\,010.2}

% latency-and-throughput.csv | system=B3_INTENT_NO_BARRIER
% step_latency_ms_median over crash-free runs only (overhead_runs_crash_free=3)
% E5 gate: runs without a declared suspend-disabled host contribute NO timing
\newcommand{\BthreeStepMedian}{2\,038.2}

% latency-and-throughput.csv | B3 median - B0 median
% = 2038.2 - 2010.2; the whole write-ahead protocol except the barrier
\newcommand{\ProtocolMinusBarrier}{28.0}

% = (2038.2 - 2010.2) / 2010.2 x 100
\newcommand{\ProtocolMinusBarrierPct}{1.4}

% latency-and-throughput.csv | AEP-full median - B3 median, both under appendfsync=everysec
% = 4004.9 - 2038.2; the two WAITAOF round trips together
\newcommand{\BarrierCost}{1\,966.7}

% half of \BarrierCost -- the protocol runs exactly two barriers per step
\newcommand{\BarrierCostEach}{983.3}

% = \BarrierCostEach / \AepStepMedian x 100
% = 983.3 / 4004.9 x 100; the end-to-end cost of adding one more barrier to the step, not the increase in the barrier bill alone
\newcommand{\ThirdBarrierStepPct}{24.6}

% latency-and-throughput.csv | \BarrierCost / \ProtocolMinusBarrier, both under appendfsync=everysec
% = (4004.9 - 2038.2) / (2038.2 - 2010.2) = 1966.7 / 28.0 = 70.14, to two significant figures
% how much more the two fsync barriers cost than everything else the protocol does
\newcommand{\BarrierToProtocolRatio}{70}

% fsync-always/analysis/latency-and-throughput.csv | AEP-full median - B3 median, both under appendfsync=always
% = 2063.4 - 2048.4
\newcommand{\BarrierCostAlways}{15.0}

% fsync-always/analysis/latency-and-throughput.csv | system=B3_INTENT_NO_BARRIER, appendfsync=always
% against 2038.2 ms under everysec: the ablated protocol's own cost is nearly policy-independent, which is what makes the two barrier figures comparable
\newcommand{\BthreeAlwaysMedian}{2\,048.4}

% fsync-always/analysis/latency-and-throughput.csv | system=AEP_FULL, appendfsync=always
% step_latency_ms_median; quoted in the threats section against the p95 below, to show the tail the three-run interval is drawn from
\newcommand{\AepAlwaysMedian}{2\,063.4}

% fsync-always/analysis/latency-and-throughput.csv | system=AEP_FULL, appendfsync=always
% step_latency_ms_p95
\newcommand{\AepAlwaysPninetyfive}{5\,067.8}

% fsync-always/analysis/latency-and-throughput.csv | system=AEP_FULL, appendfsync=always
% executions_per_second
\newcommand{\AepThroughputAlways}{0.42}

% latency-and-throughput.csv | system=AEP_FULL, appendfsync=everysec
% executions_per_second; the everysec matrix collected more runs per cell, so wall time per execution is not comparable as a benchmark -- it is quoted only against the always arm
\newcommand{\AepThroughputEverysec}{0.33}

% per-cell-metrics.csv | system=B4_DURABLE_WORKFLOW regime=(session-3) response_class=AUTHORITATIVE_READBACK
% metric=undetected_duplicate_rate | sum(successes)/sum(total) = 95/180
\newcommand{\BfourDupAuth}{0.5278}

% per-cell-metrics.csv | executions behind \BfourDupAuth
\newcommand{\BfourExecAuth}{180}

% per-cell-metrics.csv | system=B4_DURABLE_WORKFLOW regime=(session-3) response_class=POSITIVE_ONLY_READBACK
% metric=undetected_duplicate_rate | sum(successes)/sum(total) = 106/180
\newcommand{\BfourDupPosOnly}{0.5889}

% per-cell-metrics.csv | executions behind \BfourDupPosOnly
\newcommand{\BfourExecPosOnly}{180}

% per-cell-metrics.csv | system=B4_DURABLE_WORKFLOW regime=(session-3) response_class=NO_READBACK
% metric=undetected_duplicate_rate | sum(successes)/sum(total) = 101/180
\newcommand{\BfourDupNoReadback}{0.5611}

% per-cell-metrics.csv | executions behind \BfourDupNoReadback
\newcommand{\BfourExecNoReadback}{180}

% per-cell-metrics.csv | system=B4B_DURABLE_WORKFLOW_AT_MOST_ONCE regime=(session-3) response_class=AUTHORITATIVE_READBACK
% metric=lost_effect_rate | sum(successes)/sum(total) = 98/180
\newcommand{\BfourbLostAuth}{0.5444}

% per-cell-metrics.csv | executions behind \BfourbLostAuth
\newcommand{\BfourbExecAuth}{180}

% per-cell-metrics.csv | system=B4B_DURABLE_WORKFLOW_AT_MOST_ONCE regime=(session-3) response_class=POSITIVE_ONLY_READBACK
% metric=lost_effect_rate | sum(successes)/sum(total) = 91/180
\newcommand{\BfourbLostPosOnly}{0.5056}

% per-cell-metrics.csv | executions behind \BfourbLostPosOnly
\newcommand{\BfourbExecPosOnly}{180}

% per-cell-metrics.csv | system=B4B_DURABLE_WORKFLOW_AT_MOST_ONCE regime=(session-3) response_class=NO_READBACK
% metric=lost_effect_rate | sum(successes)/sum(total) = 93/180
\newcommand{\BfourbLostNoReadback}{0.5167}

% per-cell-metrics.csv | executions behind \BfourbLostNoReadback
\newcommand{\BfourbExecNoReadback}{180}

% per-cell-metrics.csv | max known_ambiguity_rate over {B4,B4b} x {all capability classes} x {all crash points}
% regime=(session-3) | 36 cells
% the durable-execution engine never declares ambiguity; that is the corner of the trilemma it cannot reach
\newcommand{\BfourFamilyAmbMax}{0.0000}

% per-cell-metrics.csv | cells behind \BfourFamilyAmbMax
\newcommand{\BfourFamilyAmbCells}{36}

% per-cell-metrics.csv | system=B3_INTENT_NO_BARRIER regime=(session-3) response_class=AUTHORITATIVE_READBACK
% metric=known_ambiguity_rate | sum(successes)/sum(total) = 0/180
\newcommand{\BthreeAmbAuth}{0.0000}

% per-cell-metrics.csv | system=B3_INTENT_NO_BARRIER regime=(session-3) response_class=POSITIVE_ONLY_READBACK
% metric=known_ambiguity_rate | sum(successes)/sum(total) = 66/180
\newcommand{\BthreeAmbPosOnly}{0.3667}

% per-cell-metrics.csv | system=B3_INTENT_NO_BARRIER regime=(session-3) response_class=NO_READBACK
% metric=known_ambiguity_rate | sum(successes)/sum(total) = 129/180
\newcommand{\BthreeAmbNoReadback}{0.7167}

% comparisons-vs-aep-full.csv | metric=undetected_duplicate_rate system=B3_INTENT_NO_BARRIER reference=AEP_FULL
% B3 0/600 vs AEP-full 0/600
% the ablation is indistinguishable from the full protocol here; that is the finding, not a caveat
\newcommand{\BthreeVsAepDupP}{1.00}

% comparisons-vs-aep-full.csv | metric=lost_effect_rate system=B3_INTENT_NO_BARRIER reference=AEP_FULL
% B3 0/600 vs AEP-full 0/600
% the ablation is indistinguishable from the full protocol here; that is the finding, not a caveat
\newcommand{\BthreeVsAepLostP}{1.00}

% comparisons-vs-aep-full.csv | metric=known_ambiguity_rate system=B3_INTENT_NO_BARRIER reference=AEP_FULL
% B3 225/600 vs AEP-full 223/600
% the ablation is indistinguishable from the full protocol here; that is the finding, not a caveat
\newcommand{\BthreeVsAepAmbP}{0.95}

% comparisons-vs-aep-full.csv | executions per arm, identical across all three metrics
\newcommand{\BthreeVsAepN}{600}

% Wilson 95% upper limit on 0/600, as a percentage
% the largest undetected-duplicate or lost-effect rate either system could have and still record zero; the difference between them is bounded by it
\newcommand{\AblationZeroUpper}{0.64}

% comparisons-vs-aep-full.csv | metric=undetected_duplicate_rate, the larger of the two arms' numerators
% B3 0/600 vs AEP-full 0/600
\newcommand{\BthreeVsAepDupCount}{0}

% comparisons-vs-aep-full.csv | metric=lost_effect_rate, the larger of the two arms' numerators
% B3 0/600 vs AEP-full 0/600
\newcommand{\BthreeVsAepLostCount}{0}

% comparisons-vs-aep-full.csv | metric=known_ambiguity_rate, |B3 successes - AEP-full successes|
% |225 - 223| over 600 executions per arm
\newcommand{\BthreeVsAepAmbDelta}{2}

% redis-kill-ablation.csv | regime=redis-kill-preack system=AEP_FULL response_class=NO_READBACK
% executions_with_an_applied_effect
\newcommand{\AepKillApplied}{10}

% redis-kill-ablation.csv | regime=redis-kill-preack system=AEP_FULL response_class=NO_READBACK
% unwanted-applied-effect rate = executions_with_an_applied_effect/executions = 10/30
% an effect put on the wire while the durability of its own intent record could no longer be confirmed
\newcommand{\AepUnwantedRate}{0.3333}

% redis-kill-ablation.csv | runs for AEP_FULL
\newcommand{\AepKillRuns}{30}

% redis-kill-ablation.csv | canary_survived/(survived+lost) for AEP_FULL
% the un-acknowledged write made immediately before the hard kill
\newcommand{\AepKillCanary}{30/30}

% redis-kill-ablation.csv | regime=redis-kill-preack system=B3_INTENT_NO_BARRIER response_class=NO_READBACK
% executions_with_an_applied_effect
\newcommand{\BthreeKillApplied}{28}

% redis-kill-ablation.csv | regime=redis-kill-preack system=B3_INTENT_NO_BARRIER response_class=NO_READBACK
% unwanted-applied-effect rate = executions_with_an_applied_effect/executions = 28/30
% an effect put on the wire while the durability of its own intent record could no longer be confirmed
\newcommand{\BthreeUnwantedRate}{0.9333}

% redis-kill-ablation.csv | runs for B3_INTENT_NO_BARRIER
\newcommand{\BthreeKillRuns}{30}

% redis-kill-ablation.csv | canary_survived/(survived+lost) for B3_INTENT_NO_BARRIER
% the un-acknowledged write made immediately before the hard kill
\newcommand{\BthreeKillCanary}{30/30}

% redis-kill-ablation.csv | B3 applied - AEP-full applied
% = 28 - 10; real non-idempotent effects that the barrier withheld under an identical injected fault
\newcommand{\UnwantedPrevented}{18}

% redis-kill-ablation.csv | Fisher exact two-tailed on [[10, 20], [28, 2]]
% AEP-full vs B3 on the unwanted-applied-effect rate
\newcommand{\UnwantedP}{1.9\times10^{-6}}

% redis-kill-ablation.csv | canary_survived + canary_lost, summed over AEP_FULL and B3_INTENT_NO_BARRIER
% un-acknowledged writes made immediately before a hard Redis kill, inside the ablation cells
\newcommand{\KillCanaryN}{60}

% experiments/results/g2-flakey-write-loss*.json | 3 replication(s) of experiments/flakey_write_loss.py
% independent runs of the probe, each rebuilding the device stack and the filesystem from scratch
\newcommand{\FlakeyReplications}{3}

% experiments/results/g2-flakey-write-loss*.json | 3 replication(s) of experiments/flakey_write_loss.py
% countable trials: the acknowledged write survived, so the trial says something about the unacknowledged one
\newcommand{\FlakeyN}{90}

% experiments/results/g2-flakey-write-loss*.json | 3 replication(s) of experiments/flakey_write_loss.py
% WAITAOF-acknowledged writes still present after the device stopped accepting writes
\newcommand{\FlakeyAckSurvived}{90/90}

% experiments/results/g2-flakey-write-loss*.json | 3 replication(s) of experiments/flakey_write_loss.py
% un-acknowledged writes destroyed by the same event
\newcommand{\FlakeyUnackLost}{90/90}

% experiments/results/g2-flakey-write-loss*.json | 3 replication(s) of experiments/flakey_write_loss.py
% narrowest write-to-write-loss exposure window, ms
\newcommand{\FlakeyWindowMin}{12.4}

% experiments/results/g2-flakey-write-loss*.json | 3 replication(s) of experiments/flakey_write_loss.py
% widest write-to-write-loss exposure window, ms; the appendfsync everysec period it sits inside is 1000 ms
\newcommand{\FlakeyWindowMax}{61.8}

% experiments/results/g2-flakey-write-loss*.json | 3 replication(s) of experiments/flakey_write_loss.py
% Fisher exact two-tailed, acknowledged vs un-acknowledged loss over the same 90 trials
\newcommand{\FlakeyBarrierP}{2.2\times10^{-53}}

% reports/raw/e1-durability-window.txt (0/10 lost under docker kill -s KILL) vs experiments/results/g2-flakey-write-loss*.json | 3 replication(s) of experiments/flakey_write_loss.py
% Fisher exact two-tailed across the two fault classes
\newcommand{\FlakeyVsProcessKillP}{5.8\times10^{-14}}

% analysis/per-execution.csv | cluster bootstrap over runs, 10000 resamples, seed 20260806, appendfsync=everysec
% 2.5th percentile of (median AEP-full - median B3); point estimate 1968.1 ms from 3 and 3 runs
\newcommand{\BarrierCostLow}{477.9}

% analysis/per-execution.csv | 97.5th percentile of the same bootstrap, appendfsync=everysec
\newcommand{\BarrierCostHigh}{1\,978.8}

% analysis/per-execution.csv | cluster bootstrap over runs, 10000 resamples, seed 20260806, appendfsync=always
% 2.5th percentile of (median AEP-full - median B3); point estimate 16.5 ms from 3 and 3 runs
\newcommand{\BarrierCostAlwaysLow}{-1\,474.6}

% analysis/per-execution.csv | 97.5th percentile of the same bootstrap, appendfsync=always
\newcommand{\BarrierCostAlwaysHigh}{22.7}

% analysis/coverage.json | runs
\newcommand{\RunsCollected}{432}

% analysis/coverage.json | executions
\newcommand{\ExecutionsCollected}{3\,780}

% analysis/coverage.json | cells
\newcommand{\CellsCollected}{126}

% analysis/coverage.json | bootstrap_resamples
\newcommand{\BootstrapResamples}{10\,000}

% reports/raw/e1-durability-window.txt | un-acknowledged writes lost to `docker kill -s KILL` under appendfsync=everysec
% everysec defers the fsync, not the write; a still-running kernel flushes the page cache
\newcommand{\ProcessKillUnackLost}{0/10}

% reports/raw/e1-durability-window.txt | usable trials
\newcommand{\ProcessKillTrials}{10}

% reports/raw/e1-durability-window.txt | min write->death over 10 trials, ms
\newcommand{\ProcessKillWindowMin}{419}

% reports/raw/e1-durability-window.txt | max write->death over 10 trials, ms
% every kill landed inside the 1000 ms fsync period, which is what makes the zero above a measurement rather than a miss
\newcommand{\ProcessKillWindowMax}{992}

% lines of Python under aep_core/, excluding __pycache__
% 14 files; regenerated on every run of this script
\newcommand{\CoreLoc}{6\,849}

% lines of Python under experiments/, excluding __pycache__
% 85 files; regenerated on every run of this script
\newcommand{\HarnessLoc}{21\,458}



% =====================================================================
% ANONYMITY TOGGLE (Monday audit, fix F2).
%
% The manuscript shipped `Anonymous Author(s)` and a URL naming a personal
% GitHub account at the same time: under double-anonymous review the URL
% defeats the anonymisation, and without it the author block is unfinished.
% Rather than pick one, the source now builds both from identical text:
%
%   public build      bash scripts/build_paper.sh              -> main.pdf
%   anonymous build   bash scripts/build_paper.sh --anonymous  -> main-anon.pdf
%
% The flag defines \ANONYMOUS before this file is read, so neither build has
% any prose of its own and the two cannot drift apart.
\newif\ifanonymous
\anonymousfalse
\ifdefined\ANONYMOUS \anonymoustrue \fi

% The artifact pointer lives here, once. Every place that cites the artifact
% uses one of these two macros, so there is exactly one block to audit for a
% leak. Two forms because the two sites need different grammar: \artifacturl
% is a bare pointer after a colon, \artifactavail completes "... are ...".
% The public build renders exactly as it did before this toggle existed.
% The author identity is gated by the SAME switch, and lives in the SAME
% block, for the same reason: one place to audit for a leak rather than two.
% The anonymous branch names nobody and adds no correspondence footnote, so
% an identity added to the wrong branch cannot reach main-anon.pdf.
% The address is set in \texttt rather than \url so it produces no link
% annotation -- an annotation is a second surface an anonymity check has to
% clear, and the Monday audit found four of them the text check had missed.
% ---------------------------------------------------------------------
% THE ARCHIVE DOI. One line, and it is the next line.
%
% Minting it is a manual web upload (docs/29-archive-deposit.md); until that
% happens this reads PENDING and every site that mentions the archive renders
% the honest "not yet deposited" sentence. Inserting the real DOI is a
% one-line edit HERE and nowhere else -- \archiveavail below derives every
% rendering from it, and no section file contains a DOI string.
% THREE states, not two. A reserved Zenodo DOI is minted but does not
% resolve until the record is published, so the deposited-at sentence
% below would be false today. \archivedoistate selects which sentence
% \archiveavail becomes; the anonymous branch reads none of them.
\newcommand{\archivedoi}{10.5281/zenodo.22766567}
\newcommand{\archivedoistate}{RESERVED}
\newcommand{\archivedoipendingmarker}{PENDING}
\newcommand{\archivedoireservedmarker}{RESERVED}
% ---------------------------------------------------------------------

\ifanonymous
  \newcommand{\artifacturl}{available via the submission system}
  \newcommand{\artifactavail}{available via the submission system}
  % A Zenodo record names its depositor, so under double-anonymous review the
  % DOI defeats the anonymisation exactly as the GitHub URL did. It gets the
  % same treatment, in the same block, for the same reason -- and note that
  % this branch does not test \archivedoi at all: an anonymous build cannot
  % leak a DOI even if one is inserted above.
  \newcommand{\archiveavail}{available via the submission system}
  \newcommand{\authorname}{Anonymous Author(s)}
  \newcommand{\authorcontact}{}
  % IEEE requires the AI disclosure in the acknowledgment of the article
  % SUBMITTED. TSE is double-anonymous, so the submitted article is this
  % build, and the requirement attaches here too. What is stripped is only
  % what identifies: no repository, no host, no script paths, no phase
  % record. What survives is the disclosure itself -- which tools, over
  % what scope, under whose direction, and who is responsible.
  \newcommand{\aidisclosure}{%
    The author used generative AI coding and drafting assistants
    throughout this work (implementation, harness, analysis pipeline,
    documentation and manuscript prose) under the author's direction
    and review. No AI system is an author; the author is solely
    responsible for the content of this article, including all claims,
    analysis and conclusions.}
  % The identifying string must not survive in the document properties either,
  % which is where an anonymised PDF usually leaks. Set empty explicitly rather
  % than relying on them happening to be unset.
  \hypersetup{pdfauthor={},pdftitle={},pdfsubject={},pdfkeywords={},
              pdfcreator={},pdfproducer={}}
\else
  \newcommand{\artifacturl}{%
    \url{https://github.com/hafizmirhamza276-lab/Research-paper-AEP}}
  \newcommand{\artifactavail}{available at%
    \ \url{https://github.com/hafizmirhamza276-lab/Research-paper-AEP}}
  \newcommand{\authorname}{Hamza Khan}
  % ORCID lives in this branch ONLY. It identifies the author as surely as
  % the name does, so under double-anonymous review it gets the same
  % treatment as the name, the e-mail and the GitHub URL.
  % IEEEtran compsoc's convention for a single-author journal paper: the
  % affiliation is a first-page \thanks footnote, not a line beside the
  % name, and it reads "A. Author is with the <unit>, <city>, <country>".
  % It sits in this branch with the name, the e-mail and the ORCID because
  % an institution identifies an author under double-anonymous review as
  % surely as the name does. The anonymous branch defines \authorcontact
  % empty, so there is no string here for that build to leak.
  \newcommand{\authorcontact}{%
    \thanks{H. Khan is with the National University of Computer and
      Emerging Sciences (FAST-NUCES), Islamabad, Pakistan.
      E-mail: \texttt{hafizmirhamza276@gmail.com}.
      ORCID: \texttt{0009-0005-9380-2188}.}}
  % The full disclosure. Third person and singular: IEEE's Editorial Style
  % Manual requires the acknowledgment be written in the third person, and
  % this is a single-author paper.
  \newcommand{\aidisclosure}{%
    The author used Anthropic's Claude (Claude Code, Opus and Sonnet
    models) and OpenAI's Codex as coding and drafting assistants
    throughout this work. Their use was not confined to one part of it:
    the protocol implementation, the fault-injection harness, the
    analysis pipeline, the verification gates, the supporting
    documentation and the prose of this manuscript were all produced
    with that assistance, under the author's direction and subject to
    the author's review. The exact model identifiers are not listed here
    because such a list goes stale; they are recorded in the artifact
    repository's commit trailers and in the per-assistant
    configuration files committed beside them, which are the
    authoritative record. The phase prompts issued from Phase~8 onward
    are committed in that repository, each before the data of the phase
    it governs;
    Phases~1A--7 carry no such record, and the project's own audit
    records that gap rather than reconstructing it.

    No AI system is an author. Every measurement reported here was
    collected by the harness on the author's host and is reproducible
    from the archived raw runs; every number in the manuscript is
    generated from those results by the artifact's table
    generator and re-derived by its independent number check. The
    author is solely responsible for the content of this article,
    including all claims, analysis and conclusions.}
  % Rendered from \archivedoi, never from a literal. While it is PENDING the
  % sentence says so; the moment a DOI replaces it the same sentence becomes a
  % true availability claim, with no other file edited.
  \ifx\archivedoistate\archivedoipendingmarker
    \newcommand{\archiveavail}{prepared and verified but not yet deposited; no
      DOI exists at the time of writing}
  \else\ifx\archivedoistate\archivedoireservedmarker
    % Reserved, not registered. The identifier is fixed and the record is
    % a draft; the DOI begins resolving only when that record is
    % published. Saying "deposited at" here would be false.
    \newcommand{\archiveavail}{prepared and verified, with the Zenodo record
      reserved under DOI \ \texttt{\archivedoi}; the identifier is fixed and
      begins resolving when the record is published}
  \else
    \newcommand{\archiveavail}{deposited at%
      \ \texttt{https://doi.org/\archivedoi}}
  \fi\fi
\fi
% =====================================================================

\begin{document}

% The break is after "Non-Idempotent", not after "APIs": at title size the
% first line holds 42 characters and not 46, so breaking after "APIs" made
% IEEEtran wrap it anyway and orphan "APIs" onto a line of its own. Verified
% in the built PDF, not assumed.
\title{AEP: Declared Ambiguity for Non-Idempotent\\
APIs Without Idempotency Keys}

\author{%
  \IEEEauthorblockN{\authorname}%
  \thanks{Manuscript prepared \today. Artifact: \artifacturl.}%
  \authorcontact%
}

\markboth{IEEE Transactions on Software Engineering}%
{AEP: Declared Ambiguity for Non-Idempotent APIs}

\IEEEtitleabstractindextext{%
\begin{abstract}
Many enterprise APIs are non-idempotent, accept no idempotency key, and cannot
be asked afterwards whether a mutation was applied. A caller that crashes around
such a call (an autonomous agent, a workflow engine) has no safe option:
retrying risks a second real-world effect nobody observes, not retrying an
effect no record accounts for. We argue this is not an engineering-quality
problem but a three-way trade whose third corner a system must make reachable:
silent failure converted into \emph{declared, durable, bounded ambiguity} an
operator can act on.

We present AEP, a fail-closed protocol in which every external side
effect is preceded by a durably acknowledged write-ahead intent, every state write is
fenced by lock ownership and an expected-version CAS in one atomic script, and
every unresolvable outcome escalates rather than guessing.
We evaluate it against five baselines under \texttt{SIGKILL} faults,
across endpoint capabilities differing in what a read-back can settle.

Our central result separates two claims usually sold as one.
\emph{Detection} (no undetected duplicate and no lost effect, leaving a
residual of declared ambiguity) comes from the pre-dispatch record and a
transition table that forbids re-entry into dispatch, not from the durability
barrier. \emph{Prevention} is the barrier's, against a narrower
fault than it appears: it withholds dispatch when the store dies inside the
acknowledgment window, not when storage discards writes while reporting
success, which we injected and where it dispatched. Because detection
does not depend on the barrier for record-preserving faults, its cost is a
deployment choice, not the protocol's price.
\end{abstract}

\begin{IEEEkeywords}
Fault tolerance, non-idempotent APIs, idempotence, distributed coordination,
write-ahead logging, fault injection, reliability engineering, autonomous
agents.
\end{IEEEkeywords}}

\maketitle
\IEEEdisplaynontitleabstractindextext
\IEEEpeerreviewmaketitle

\section{Introduction}
\label{sec:introduction}

An autonomous agent is asked to post a journal entry, issue a refund, or file a
regulatory notification. It calls an enterprise API written years before any of
this. The call times out, and nothing on either side can say whether the money
moved. Autonomous agents are the setting in which this is now most visible, and
we use them throughout as the motivating example. The problem itself belongs to
the endpoint, and any caller that can fail mid-call inherits it. A workflow
engine, a retrying batch job and a human operator with a script are all in the
same position.

This is an old problem in a new setting, and the new setting removes the
assumption that the old solutions rest on. Durable-execution engines, workflow
systems and message brokers achieve their guarantees by requiring that the
operation being retried is idempotent, or that the remote side will accept a
deduplication key, or that both endpoints can be enrolled in one
transaction~\cite{helland2007apostate,burckhardt2021durable,zhuang2023exoflow}.
Temporal's documentation is explicit about the precondition: ``You should
always make your business logic Activities idempotent in Temporal. Because
Activities may be retried, these functions may be executed more than
once''~\cite{temporal-activity-definition}. The legacy endpoints at issue satisfy
none of these conditions. They are non-idempotent and they accept no idempotency
key. Many of them also cannot be asked afterwards whether a mutation was
applied, and that third gap is the one that does the real damage.

\subsection{The trilemma}
\label{sec:trilemma}

Consider a worker that has sent a mutation and then crashed, and a recovery
process that must decide what to do. Exactly three things can happen to the
external world, and a system must choose which one it is willing to produce.

\begin{itemize}
\item An \textbf{\undetdup{}}. The recovery process re-sends. If the first
      request had in fact been applied, a second real effect now exists and no
      component of the system knows it. A naive retry produces this outcome,
      and so, as \cref{sec:evaluation} shows, do a lease, a fenced
      compare-and-swap and a durable event-sourced history, which each fail
      to prevent it because none of them is a record written before the
      call.
\item A \textbf{\lostfx{}}. The recovery process does not re-send and records
      the attempt as failed. If the first request was applied, a real effect
      now exists that the system's records deny. An at-most-once configuration
      produces this outcome. It is quieter than a duplicate and no better: a
      refund that the ledger says did not happen becomes a reconciliation
      problem that surfaces in the accounts weeks later.
\item A \textbf{\declamb{}}. The system declines to decide, records that it
      cannot decide, and stops. The effect may or may not exist. What the
      system guarantees is that a durable record says so, and that no further
      automated action is taken on that execution.
\end{itemize}

The first two outcomes are silent and the third is not. Our position is that
the third corner is the one a protocol should be engineered to reach. Against
an endpoint that cannot be queried, avoiding ambiguity altogether is
impossible, so the quantity a system should be judged on is whether its
residual uncertainty is declared and bounded rather than silent and unbounded.

This impossibility is not incidental. Deciding whether a remote effect occurred,
when the channel may fail and the remote side may not answer, is the
coordinated-attack problem; no protocol resolves it with certainty, and the
classical results on consensus under crash faults and on failure detection say
why~\cite{fischer1985impossibility,chandra1996failuredetectors}. What a designer
can choose is where the uncertainty is allowed to surface. We call our protocol
AEP. It records an intent before every external call and fails closed when
recovery cannot resolve the outcome. AEP places the uncertainty in a durable
state that an operator can see, and in every cell we measured, none of it
reached the accounts.

\subsection{Contributions}

\begin{itemize}
\item[\textbf{C1}] \textbf{The declared-ambiguity formulation.} We state the
      problem as the three-way trade of \cref{sec:trilemma} rather than as a
      pursuit of exactly-once. The formulation gives three properties: fenced
      state, detectable ambiguity, and a fail-closed liveness bound. Each is
      mapped to the code path that enforces it, and each has its residual
      window declared rather than assumed away (\cref{sec:model,sec:protocol}).
\item[\textbf{C2}] \textbf{A protocol and an implementation} in which a
      durably acknowledged write-ahead intent is a checked precondition of
      dispatch authority, not a matter of call ordering. Under a trusted-code
      assumption, the fsync acknowledgment returns an opaque, non-copyable,
      single-use, scope-bound in-process dispatch guard. The guard is consumed
      to write a Redis-visible authorization, and the pre-dispatch script
      refuses to proceed without it (\cref{sec:protocol}). This is a
      control-flow invariant of the supported API. It is not a cryptographic
      capability, and it is not a boundary against arbitrary code in the same
      Python process.
\item[\textbf{C3}] \textbf{An evaluation under real process kills} across six
      named crash points, three endpoint reconciliation capabilities and
      five reported fault regimes, against five baseline designs, one of them
      in two configurations. The five regimes are every execution crashed,
      none crashed, Redis hard-killed before the acknowledgment, Redis killed
      after dispatch, and storage discarding writes while reporting success.
      (A sixth regime, a $30\%$ crash probability, was collected; its rates
      are not in these tables and the collection is in the artifact.)
      % Regimes counted from each collection root's analysis/coverage.json.
      % The five reported: crashed and p0 from experiments/results/matrix,
      % redis-kill-preack from matrix and the b2/phase13-armA sessions,
      % redis-kill-inflight from phase13-inflight-s1,s2 (6.3.2.5), and
      % write-loss-preack from reports/raw/ws4-writeloss-s1-2026-09-07
      % (6.3.4). The sixth is p30, ws5-2026-09-10/t2-p30, which no table
      % quotes. (session-3) and crashed are one regime under the two labels
      % docs/29 3 declares as a normalisation, so they are counted once.
      % This said 'three collected fault regimes' until 2026-09-17, which
      % was the matrix root's three and excluded the two regimes 6.3.2.5
      % and 6.3.4 report -- one of them the fault the abstract's prevention
      % sentence turns on.
      AEP records no
      undetected duplicate and no lost effect in any cell measured, and the
      baselines without a pre-dispatch record duplicate in most crashed
      executions. AEP's residual is declared ambiguity, whose rate is a
      function of what the endpoint can be asked (\cref{sec:evaluation}).
\item[\textbf{C4}] \textbf{A decomposition of the mechanism, by ablation,
      that reassigns our own headline result.} The write-ahead pattern is
      normally presented as one mechanism delivering one guarantee, and it is
      two. The pre-dispatch record together with no re-entry
      is what makes an outcome detectable, and it does so without the
      durability barrier: ablating the barrier produces no observed difference
      in the crashed-regime detection metrics (\cref{sec:eval-detection}). The
      barrier contributes something else, withholding the external call when
      durability can no longer be confirmed. That contribution is invisible to
      the crash faults that dominate this literature and large under the fault
      that targets Redis (\cref{sec:eval-prevention}). Separating the two
      changes what a deployment must pay for: detection is nearly free,
      prevention is where the fsync cost lives, and an operator can buy the
      first without the second (the supplementary material, \emph{The barrier
      is a deployment choice}).
\end{itemize}

\subsection{What this paper does not claim}

We claim no exactly-once external execution, no prevention of a duplicate the
provider has already accepted, no high availability, consensus or split-brain
immunity, and no durability beyond one local AOF fsync. The system is a single
Redis trust domain and the measurements are single-node. These are stated as a
table of non-claims in \cref{sec:model} with the reason each is out of scope,
and revisited in \cref{sec:threats}.

\section{Motivating Traces}
\label{sec:motivating}

The three traces below are not hypotheticals. Each is a run of our harness
(\cref{sec:evaluation}), replayed from its event log and cross-checked against
the mock provider's ground-truth ledger --- an independently written record of
every mutation the provider \emph{actually applied}. In each, the agent
workload is the same: acquire work, call one non-idempotent endpoint, record
the outcome.

\subsection{Trace 1 --- the duplicate nobody sees (B0, naive retry)}

A worker sends the mutation. The provider applies it. The worker is
\texttt{SIGKILL}ed after the response is on the wire but before the outcome is
recorded (crash point \texttt{after\_response\_before\_resolution}). A
supervisor observes an execution with no terminal record and runs the step
again. The provider applies the mutation a second time.

The ground-truth ledger now holds two applied mutations for one logical
step. The system's own records hold one successful execution. Nothing in the
system is in a state that a monitor could alert on: there is no error, no
retry counter above its threshold, no ambiguous record. The execution is
\emph{green}.

This is not a strawman configuration; it is what a retry-on-timeout loop does,
and it is the behaviour we measure for B0. Across the crashed cells B0 ends
this way in \BaselineDupLow{}--\BaselineDupHigh{} of executions depending on
the endpoint, and so do B1 and B2 (\cref{tab:outcomes}).
% \BaselineDupLow / \BaselineDupHigh: generated/numbers.tex, the min and max
% undetected-duplicate rate over {B0,B1,B2} x three response classes,
% crashed regime.

\subsection{Trace 2 --- a lease and a fenced write do not help}

The same crash, against B1 (a real Redis lease) and B2 (a lease plus an exact
expected-version compare-and-swap on the state document). Both are the patterns
an experienced engineer reaches for, and both are doing their job: no stale
writer corrupts the state document, and no two workers hold the lock at once.

Neither prevents the second effect. The lease governs \emph{who may write
state}; the CAS governs \emph{which version of state may be replaced}. The
external call is neither. When the supervisor re-runs the step it holds a
valid, freshly acquired lease and writes a correctly versioned record --- of a
duplicate. \Cref{sec:evaluation} shows B1 and B2 duplicating at
substantially the same rate as B0, which is the point: the mechanisms that
protect the database do not protect the world.

\subsection{Trace 3 --- a durable log is necessary and not sufficient (B4)}

B4 is an event-sourced durable-execution engine of the kind these systems are
built on: an append-only history in Redis, each append acknowledged durable
with the \emph{same} \texttt{WAITAOF} barrier AEP uses, replay after a crash,
and completed activities memoised so they are not re-run. It is not ablated on
durability. Its history is on disk before the call.

It duplicates anyway. The history records that the activity was
\emph{scheduled}; it does not, and cannot, record whether the provider applied
the effect, because that fact was never available to the engine. On replay the
engine sees a scheduled activity with no completion and --- following the
vendor's documented default, where Activity Maximum Attempts is
unlimited~\cite{temporal-retry} --- schedules it again.

The mitigation the vendor recommends is to make the activity
idempotent~\cite{temporal-activity-definition}. That is exactly the
precondition this paper's premise removes.

\subsection{Trace 3b --- and the at-most-once configuration loses the effect}

Setting Maximum Attempts to 1 is documented and standard: ``Setting the value
to 1 means a single execution attempt and no retries''~\cite{temporal-retry}.
We run it as B4b. It cannot produce a caller-caused duplicate, and it does not.

What it produces instead is the other silent failure. Where the provider
applied the mutation before the worker died, B4b's history records an activity
that timed out; the workflow fails; the effect exists and the records deny
it. Critically, that record is an ordinary failure and not a declared
ambiguity --- there is no outcome class in B4b in which the engine can say
\emph{I do not know}, so nothing escalates and no operator is told that a
real-world effect may be unaccounted for.

\subsection{What the traces establish}

\begin{table}[t]
\centering
\caption{The trilemma, as the systems under comparison instantiate it. A
durable write-ahead record is necessary (Traces 1, 2) and not sufficient
(Trace 3); what distinguishes AEP is an outcome class in which it can decline
to decide. B3 --- AEP with the durability barrier removed and nothing else ---
reaches the same corner, which is the point: the barrier is not what buys
declared ambiguity (\cref{sec:eval-detection} measures the ablation), and
\cref{sec:eval-prevention} isolates what it does buy.}
\label{tab:trilemma}
\small
\begin{tabular}{@{}lccc@{}}
\toprule
 & \undetdup{} & \lostfx{} & \declamb{} \\
\midrule
B0 naive retry            & \textbf{yes} & --- & no \\
B1 lease-only             & \textbf{yes} & --- & no \\
B2 CAS-only               & \textbf{yes} & --- & no \\
B4 durable, max attempts $\infty$  & \textbf{yes} & --- & no \\
B4b durable, max attempts $=1$     & no & \textbf{yes} & no \\
\midrule
B3 AEP minus the barrier  & no & no & \textbf{yes, bounded} \\
AEP-full                  & no & no & \textbf{yes, bounded} \\
\bottomrule
\end{tabular}
\end{table}

\Cref{tab:trilemma} is the shape of the argument. At-most-once dispatch is
available off the shelf --- B4b buys it with one configuration setting, and so
does AEP's ablation B3. What is not available off the shelf is knowing
\emph{which} of duplicate-or-loss occurred, because the fact a system would
need for that --- whether the provider applied the effect --- is not in any log
it can write and cannot be put there by any amount of logging. It has to be
obtained from the endpoint, and \cref{sec:eval-rq1} shows the price of that: the
achievable outcome is bounded by what the endpoint can be asked.

\section{System and Failure Model}
\label{sec:model}

This section condenses a formal model held in the
artifact, which states every claim below with a citation to the enforcing
code by file and line and is machine-checked in CI: a script fails the build if a
cited path stops existing or a cited line falls outside its file. That check
proves the citation ranges are valid, not that the semantics are right. The
semantic evidence is the per-property regression tests named in
\cref{sec:protocol}.

\subsection{Principals and store}

A \emph{runner} performs at most one external mutation per invocation. A
\emph{recovery service} is read-only with respect to the external system and
writes only coordination state. A \emph{connector} is the only component that
touches the external API. The store is a single self-hosted Redis~7.2 instance
with the append-only file enabled and RDB snapshots disabled. There is no
replica, no Sentinel and no Cluster.

Redis executes a Lua script without interleaving, so a script is a single
serialization point. This is the only atomicity primitive the protocol has and
every invariant below is expressed inside one.

\subsection{Clocks}
\label{sec:model-clocks}

There is no synchrony assumption beyond lease TTLs. All ordering decisions
between principals are taken against the Redis server clock; no cross-worker
clock comparison occurs. Process-local monotonic time is used only for
heartbeat intervals and watchdogs, never for an authoritative decision. The
lease budget $T_{\text{client}} \le T_{\text{lock}} - B$ with $B \ge 15$\,s is
enforced at configuration time. It is a timing budget rather than a
guarantee: a scheduling gap remains between a successful preflight and
transmission, and we declare it as such.

\subsection{The external endpoint}
\label{sec:model-endpoint}

The endpoint is non-idempotent and non-cooperative: it need not accept an
idempotency key, and its response may be absent, late, or self-contradictory.
Two things about it are modeled explicitly, and the second is the axis the
evaluation turns on.

\paragraph{Evidence at mutation time} Ambiguity is the default here, not a
special case. The runner accepts only policy-declared
success and failure values, required non-empty and disjoint; any other value,
and any exception whatever, is coerced to ambiguity.

\paragraph{Reconciliation capability} What the endpoint can be asked after the
fact is a declared, typed property with three values
(\cref{tab:capabilities}). The declaration fails closed: a missing, null or
unrecognized capability is treated as the weakest, not as the strongest.
Classification is total. In particular, ``the endpoint reports the mutation
absent'' combined with a positive-only capability is a \emph{named contract
violation} that resolves to ambiguity, not a fall-through to ``it did not
happen''.

\begin{table}[t]
\centering
\caption{Endpoint reconciliation capability. The right-hand column is the
ceiling on what any protocol can conclude, and \cref{sec:eval-rq1} shows
AEP's declared-ambiguity rate tracking it.}
\label{tab:capabilities}
\small
\begin{tabular}{@{}p{0.30\columnwidth}p{0.60\columnwidth}@{}}
\toprule
Capability & Conclusions the endpoint can support \\
\midrule
\textsc{authoritative} \textsc{readback} &
  applied, not-applied, unknown, conflict. Absence \emph{proves} the mutation
  did not occur. \\
\textsc{positive-only} \textsc{readback} &
  applied, unknown, conflict. Presence can be confirmed; absence proves
  nothing. \\
\textsc{no-readback} &
  none. The endpoint cannot be queried at all. \\
\bottomrule
\end{tabular}
\end{table}

\subsection{Failure model}

The model is crash-and-delay, not Byzantine. Redis is trusted to execute
scripts atomically
and honor expiry; the provider is trusted only insofar as its declared
evidence is truthful.

\begin{itemize}
\item[\textbf{F1}] \textbf{Worker crash at any instruction boundary.} Six named
      crash points around the external call
      (\cref{tab:crashpoints}). In the evaluation these are OS-level
      \texttt{SIGKILL} of a separate worker process, not in-process
      exceptions.
\item[\textbf{F2}] \textbf{Partition between a worker and Redis.} Any Redis
      command may raise; every raise forbids progress rather than permitting a
      guess. A partition outlasting the lease is observationally equal to F5.
\item[\textbf{F3}] \textbf{Redis restart with AOF replay.} Up to roughly one
      second of acknowledged-but-unsynced writes may be absent under
      \texttt{appendfsync everysec}~\cite{redis-persistence}.
      \Cref{sec:eval-durability} reports two probes that sharpen this
      considerably. The set of faults that can actually realize that loss is
      narrower than the documentation's wording suggests, and a process kill
      is not among them; a host-level write loss is, and realizes the loss
      every time.
\item[\textbf{F4}] \textbf{Delayed or duplicated external responses.} A late
      effect cannot be recalled; recovery waits out a reconciliation delay
      before drawing any conclusion, and never re-dispatches.
\item[\textbf{F5}] \textbf{Worker pause past lease expiry} (GC or VM stall),
      handled by an atomic pre-dispatch preflight, cancellation of the
      protected task on ownership loss, and CAS rejection of the late write.
\end{itemize}

Out of model: Byzantine Redis; a provider whose declared evidence contradicts
its own effects; loss of the Redis host once the fsync acknowledgment has been
returned; adversarial writers with direct socket access; compromise of vault
keys.

\begin{table}[t]
\centering
\caption{The six crash points, and what exists in the world at each. The
right-hand column is why the declared-ambiguity rate in \cref{sec:eval-rq1} varies
so sharply across them.}
\label{tab:crashpoints}
\small
\renewcommand{\arraystretch}{1.08}
\begin{tabularx}{\columnwidth}{@{}>{\raggedright\arraybackslash}p{0.46\columnwidth}>{\raggedright\arraybackslash}X@{}}
\toprule
Crash point & State when the worker dies \\
\midrule
\path{before_intent_write}            & no record, no effect \\
\path{after_intent_before_barrier}     & record in memory, maybe not on disk; no effect \\
\path{after_barrier_before_dispatch}   & record durable; \textbf{no effect} \\
\path{mid_dispatch}                    & record durable; effect \textbf{likely} \\
\path{after_response_before_resolution} & record durable; effect known to the dead process only \\
\path{after_resolution_before_barrier} & outcome written, not yet acknowledged \\
\bottomrule
\end{tabularx}
\end{table}

\subsection{Non-claims}
\label{sec:nonclaims}

We state these because reviewers are right to look for them, and because
several are structural rather than incidental.

\begin{table}[t]
\centering
\caption{Non-claims, with the reason each is out of scope.}
\label{tab:nonclaims}
\small
\begin{tabular}{@{}p{0.36\columnwidth}p{0.54\columnwidth}@{}}
\toprule
Non-claim & Why \\
\midrule
Exactly-once external execution
  & Redis and a legacy provider share no transaction coordinator. The
    protocol converts this into declared ambiguity. \\
Duplicate \emph{prevention}
  & AEP prevents \emph{automatic} duplicates. An effect the provider already
    accepted cannot be recalled. \\
HA, consensus, split-brain immunity
  & Single instance; the lock is a lease~\cite{gray1989leases}. This is also
    why the AOF-rewind residual below has no local fix. \\
A hardened trust domain
  & No Redis ACL; any principal reaching the socket can write the state key
    directly. \\
Durability beyond one local AOF fsync
  & The barrier requires one \emph{local} fsync~\cite{redis-waitaof}; it is
    explicitly not a substitute for replication. \\
Operator alerting
  & Not implemented. ``Escalates'' means \emph{reaches a terminal automated
    state and pauses the execution}, nothing more. \\
A verified implementation
  & The properties are model-checked (\cref{sec:modelchecking}); the Python is
    not. Refinement from the specification to the code is neither proven nor
    claimed. \\
\bottomrule
\end{tabular}
\end{table}

\section{The AEP Protocol}
\label{sec:protocol}

\subsection{Ordering}

The runner's sequence, with the property each step exists to serve:

\begin{enumerate}
\item Acquire the lease with jittered retry.
\item Require that execution state already exists.
\item Prepare, vault and read back the request \emph{before any intent
      exists}, so that the intent, once written, names a request that is
      already durable.
\item Write the intent \textsc{none}$\rightarrow$\textsc{about-to-fire} by
      fenced CAS, \textbf{and} run the durability barrier on the \emph{same
      pinned connection}.
\item Convert the barrier's acknowledgement into a Redis-visible dispatch
      authorization.
\item Run the atomic pre-dispatch preflight, which re-checks that
      authorization.
\item Mint a single-use dispatch capability.
\item Only then call the connector.
\end{enumerate}

Two of these steps can refuse, and both refusals happen \emph{before} any
provider bytes exist, so both may soundly record a confirmed failure. If the
barrier fails at step~4, the runner makes one fenced attempt to record
\textsc{failed-confirmed} with evidence \texttt{LOCAL\_NO\_DISPATCH}. If the
preflight rejects at step~6 --- a lost lease, a moved version, an absent
authorization --- the runner does the same, with a reason naming the preflight
rather than the barrier, provided it still owns the lease. Either way, if
ownership has already been lost the runner records nothing and leaves the
intent for recovery. Together these are the paths that produce the measured
result of \cref{sec:eval-prevention}; the outcome classifier keys on the
resulting status rather than on which of the two refused, so they are counted
alike.

\subsection{P1 --- Fenced state}
\label{sec:p1}

\begin{property}[Fenced state]
No committed state write can be superseded and then resurrected by a stale
writer. Every write to the state key is admitted only if, in one atomic Lua
invocation, (a)~the caller's ownership token equals the live value of the lock
key, and (b)~the stored version equals the caller's expected version and the
candidate's version equals expected${}+1$.
\end{property}

Both conjuncts are necessary and the failure of either is a known
pattern. Version-only CAS admits a writer that lost its lease but holds the
current version; token-only checking admits a writer rebasing on a stale
read. The pairing is the fencing-token argument made against lease-only
locking~\cite{kleppmann2016locking}, with the token check and the write placed
inside one script so they cannot interleave. The same script additionally
freezes, inside the atomic region, the execution identity, the schema version,
every envelope field outside a declared mutable set, every other ledger entry,
and the intent's immutable fields; transitions are append-only with a prefix
check, and ledger entries can never be deleted.

\paragraph{Declared residuals} A plain \texttt{SET} from any client with socket
access is unconstrained (there is no ACL). Expiry, not deletion, ends the
guarantee: after the retention floor elapses the record is gone and any writer
may create version~1 again. And --- confirmed by a probe rather than argued ---
an AOF rewind can un-fence: P1 is a statement about one monotonic Redis
timeline, and because lock keys are deliberately \emph{not} barriered, a lease
whose release was lost can reappear alongside a rewound version, so a
previously fenced writer satisfies both conjuncts. There is no local fix; the
fix is consensus, which is a non-claim (\cref{tab:nonclaims}).

\subsection{P2 --- Detectable ambiguity}
\label{sec:p2}

\begin{property}[Detectable ambiguity]
If a \emph{durable} \textsc{about-to-fire} intent exists, then under F1--F5 it
eventually reaches exactly one \emph{terminal automated} state ---
\textsc{fired-confirmed}, \textsc{failed-confirmed} or
\textsc{permanently-ambiguous} --- and no automated transition moves it
afterwards; it is never silently re-dispatched and never silently dropped.
This holds provided the recovery loop runs, survives its own pass, can
eventually obtain the lease, and is not made eligible before the dispatch
window closes.
\end{property}

\emph{Terminal automated} is exact rather than cautious. The only edges out of
\textsc{permanently-ambiguous} are the two operator transitions of
\cref{sec:p3}, so an intent that has been escalated can still be resolved
later --- by a person, holding the lease and the exact version, asserting
something the protocol cannot check.

The design question P2 raises is what makes ``barrier before dispatch'' more
than a convention. In this protocol it is a checked precondition, in four
links:

\begin{enumerate}
\item Through the supported API, the barrier issues a \texttt{DurabilityAck}
      only when \texttt{WAITAOF} returned at least one local fsync.
\item Ordinary callers cannot construct, subclass or copy
      \texttt{DurabilityAck}; the in-process registry detects reuse and binds
      the guard to \texttt{execution:intent:prepared\_version}, so a guard for
      one attempt cannot authorise another through the supported path.
\item Authorizing dispatch \emph{consumes} the ack, then runs a script that
      re-checks lease, version, status and binding digest and writes an
      authorization key.
\item The preflight script requires that exact value; an absent authorization
      defaults to a value that cannot match, so it fails closed.
\end{enumerate}

\paragraph{The irreducible gap, stated} Redis exposes no way for a Lua script
to verify that a previous \texttt{WAITAOF} succeeded. A complete in-store proof
is therefore impossible. Under trusted-code assumptions, an authorization
records that the supported path received and consumed its ordinary barrier
guard; it is not an independent proof that Redis fsynced. The HMAC and registry
are mutation/reuse checks, not a cryptographic security boundary. Python
underscore naming and closure placement do not prevent arbitrary same-process
code from calling the module-internal issuer without running the barrier. A
principal with direct Redis access can also write the authorization key, and a
caller inside the process can compose the binding service with a connector
directly, bypassing the runner. A stronger attacker model requires a separately
enforced process or service boundary, which this implementation does not have.

\paragraph{A fifth link, and it is a clock} The four links above run from the
barrier to the dispatch. They are silent about the other end, and P2 needs it:
recovery must not classify an intent whose worker may still be about to
transmit. Nothing structural prevents that --- both principals are acting on
the same fenced state, and a lease can expire at any moment. What prevents it
is a delay. An intent becomes eligible for recovery only after
$T_{\text{client}} + B + T_{\text{settle}}$ has elapsed since it was prepared,
which is sized to cover the whole dispatch window. Remove that delay and the
protocol loses an effect: recovery reads back \emph{not applied}, truthfully,
because the request has not been sent yet, and records a confirmed failure that
the dispatch then contradicts. \Cref{sec:modelchecking} produces this in ten
steps. The delay is therefore a precondition of P2's soundness rather than a
courtesy about late responses, and it is a real-time argument --- which is why
the residual named in \cref{sec:model-clocks} is the gap it cannot close.

\paragraph{Recovery classification} For an eligible intent, under a freshly
acquired lease and after re-validating status and version, the mapping from
observation to outcome is total (\cref{tab:classification}). Recovery has no
reference to the mutation call at all --- it can only read back --- and the
transition table has no edge back to \textsc{about-to-fire}, checked in Python
and re-checked inside the atomic script.

\begin{table}[t]
\centering
\caption{Recovery classification. The two rows in bold are where declared
ambiguity is produced by design rather than by failure, and they are what
\cref{sec:eval-rq1} measures.}
\label{tab:classification}
\small
\begin{tabular}{@{}p{0.50\columnwidth}p{0.40\columnwidth}@{}}
\toprule
Observation & Outcome \\
\midrule
capability undeclared / unrecognised & \textbf{ambiguous}, no query \\
\textsc{no-readback}                 & \textbf{ambiguous}, no query \\
applied (either querying class)      & confirmed \\
not-applied $+$ \textsc{authoritative} & refuted \\
not-applied $+$ \textsc{positive-only} & \textbf{ambiguous} (named violation) \\
conflict                             & \textbf{ambiguous} \\
unknown, budget remaining            & stay unresolved, backoff \\
unknown, budget exhausted            & \textbf{ambiguous} \\
\bottomrule
\end{tabular}
\end{table}

\paragraph{False-positive ambiguity is the price, and it is the right
direction} A crash after a durable intent but \emph{before} transmission yields
an intent that recovery must treat as ambiguous even though no effect
exists. This inflates the declared-ambiguity rate; it can never inflate the
undetected-duplicate rate. \Cref{sec:eval-rq1} shows this asymmetry directly in
the data: ambiguity is highest at exactly the crash point where no effect can
possibly exist.

\subsection{P3 --- Fail-closed liveness bound}
\label{sec:p3}

\begin{property}[Fail-closed bound]
Automated reconciliation of one intent terminates within a configurable attempt
budget and duration budget, after which the intent is placed in a terminal
automated state and the execution is withdrawn from normal scheduling.
\end{property}

Defaults are 8 attempts and 24 hours, with full-jitter backoff. Attempt
counting is monotone and re-checked inside the atomic script: an
unresolved-to-unresolved edge must increment the counter by exactly one. On
exhaustion the intent becomes \textsc{permanently-ambiguous}, the execution is
paused, no new intent may be created for it, and the only exit is an operator
transition that itself requires the lease and the exact version.

\paragraph{Declared residuals} The bound is consulted only on the
\emph{unknown} branch --- every other observation resolves immediately, so this
is not a gap in practice, but it is not a global watchdog either. Lease
contention does not consume budget. There is no escalation \emph{mechanism}:
``ejects to operator escalation'' means reaching a terminal state and pausing,
and nothing alerts. And escalated records still expire: a retention floor of 31
days is enforced for any ledger containing an unresolved or ambiguous record,
but a floor is not permanence, so after it elapses the escalation evidence is
deleted without an operator having acted. Both residuals are confirmed by
executable probes in the artifact, not asserted.

\subsection{The state machine}

\Cref{fig:statemachine} is the transition set exported from the implementation,
not a drawing of it, so the presence or absence of an edge is a fact about the
code. Three states are terminal for automated reconciliation:
\textsc{fired-confirmed}, \textsc{failed-confirmed} and
\textsc{permanently-ambiguous}, which is the declared-ambiguity corner. Nothing
returns to \textsc{about-to-fire}: once an intent leaves the dispatch window it
cannot re-enter it, and that absence is P2's ``never silently re-dispatched''.
The dashed edges are the operator transitions of \cref{sec:p3}, and they are the
only way out of that corner.

\begin{figure}[t]
\centering
% GENERATED by scripts/gen_state_machine.py -- do not edit.
% Edge set imported from aep_core.core.intents
% .LEGAL_INTENT_TRANSITIONS, so this figure cannot drift from
% the transition table the Lua script enforces.
% 10 edges, 6 states, 3 terminal for automated reconciliation.
\begin{tikzpicture}[
    every node/.style={font=\scriptsize},
    state/.style={draw,rounded corners=2pt,inner xsep=4pt,inner ysep=3pt,align=center,fill=black!3},
    terminal/.style={state,fill=black!8,very thick},
    elabel/.style={font=\scriptsize,inner sep=1pt,align=center},
    to/.style={-{Stealth[length=4pt]},shorten >=1pt,shorten <=1pt},
  ]
  \node[state] (none) at (0.0,0.0) {\textsc{none}};
  \node[state] (atf) at (0.0,-1.5) {\textsc{about-to-fire}};
  \node[state] (fu) at (0.0,-3.3) {\textsc{fired-}\\\textsc{unconfirmed}};
  \node[terminal] (fc) at (-3.08,-5.6) {\textsc{fired-}\\\textsc{confirmed}};
  \node[terminal] (flc) at (3.08,-5.6) {\textsc{failed-}\\\textsc{confirmed}};
  \node[terminal] (pa) at (0.0,-5.6) {\textsc{permanently-}\\\textsc{ambiguous}};

  \draw[to] (none) edge node[elabel,midway,right=2pt] {CAS $+$ barrier} (atf);
  \draw[to,out=180,in=90,out looseness=1.5,in looseness=0.45] (atf) edge node[elabel,pos=0.46,sloped,above=1pt] {declared success} (fc);
  \draw[to,out=0,in=90,out looseness=1.5,in looseness=0.45] (atf) edge node[elabel,pos=0.46,sloped,above=1pt] {declared failure} (flc);
  \draw[to] (atf) edge node[elabel,midway,right=2pt] {anything else} (fu);
  \draw[to,out=200,in=160,looseness=2.8] (fu) edge node[elabel,midway,left=2pt,align=right] {unknown;\\budget left} (fu);
  \draw[to,out=225,in=45] (fu) edge node[elabel,pos=0.63,sloped,above=2pt] {read-back applied} (fc);
  \draw[to,out=315,in=135] (fu) edge node[elabel,pos=0.58,sloped,above=2pt] {authoritative absence} (flc);
  \draw[to] (fu) edge node[elabel,pos=0.72,right=2pt] {budget spent} (pa);
  \draw[to,dashed,out=170,in=10] (pa) edge (fc);
  \draw[to,dashed,out=10,in=170] (pa) edge (flc);
\end{tikzpicture}

\caption{The intent state machine, generated from the transition table in
\texttt{aep\_core/core/intents.py}. Solid edges are automated; the dashed edges
from \textsc{ambiguous} are the operator-only transitions. There is no edge
back into \textsc{about-to-fire} --- that absence is what ``never silently
re-dispatched'' means, and it is enforced inside the atomic script rather than
by the caller.}
\label{fig:statemachine}
\end{figure}

\subsection{Model checking}
\label{sec:modelchecking}

We specify the protocol in TLA+ and check it with TLC. The specification is
transcribed from the implementation --- the Lua guards, the exhaustive
transition table, the creation fence, the recovery classification --- and
deliberately not from this section. A gate re-derives the transition table and
the status alphabet from the source and fails if the two disagree, because a
model checker asked about the wrong state machine does not complain: it reports
success.

The endpoint is modelled as \cref{sec:model-endpoint} describes it, and this is
the modelling decision the result rests on. It is never given an intent
identifier, so it cannot deduplicate and a second dispatch yields a second
effect; whether an effect is applied is chosen adversarially rather than
derived from anything the caller did; and ambiguous evidence is reachable on
every path, so no caller can ever separate \emph{no effect} from \emph{effect,
answer lost}. A god-view variable records what really happened at the endpoint,
and no protocol action may read it --- enforced textually across the
specification's protocol definitions, because a guard that could see it would
verify a protocol strictly stronger than the implementation while every
property still came back green.

\emph{Checked:} P1 as version monotonicity and the absence of status
regression; P2 as legal-edges-only, no re-entry into \textsc{about-to-fire}, at
most one unresolved intent, and eventual resolution under fairness; P3 as the
bounded attempt budget; and the trilemma as two invariants --- no undetected
duplicate, and no lost effect. Declared ambiguity is not a violation, and that
asymmetry is \cref{sec:trilemma} stated formally.

\begin{table}[t]
\centering
\caption{Seven of the nine configurations that must fail. Each removes one
stated assumption; \emph{steps} is the length of the counterexample TLC
returns. Two further configurations, not shown, are non-vacuity witnesses:
\textsc{p2}'s liveness property is an implication, and they establish that its
antecedent is reachable rather than letting it pass while meaning nothing.}
\label{tab:modelchecking}
\small
\begin{tabular}{@{}p{0.44\columnwidth}p{0.34\columnwidth}r@{}}
\toprule
Assumption removed & Property that breaks & Steps \\
\midrule
% Steps are counterexample depths from scripts/run_tlc.sh at the committed
% bounds; one row per file in formal/configs/. Depths are used rather than
% state counts because a failing run stops at its first counterexample, so its
% state count reflects the search order and is not reproducible.
A single Redis timeline   & P1 monotonicity      & 4 \\  % aof-rewind.cfg
A truthful \texttt{fsync} & no lost effect       & 8 \\  % write-loss.cfg
The barrier, plus a restart & no lost effect     & 8 \\  % b3-no-barrier-restart.cfg
A truthful endpoint       & no lost effect       & 9 \\  % untruthful-endpoint.cfg
\quad\emph{same, duplicate horn} & no undetected duplicate & 16 \\ % untruthful-endpoint-duplicate.cfg
The reconciliation delay  & no lost effect       & 10 \\ % no-reconcile-delay.cfg
Operator authentication   & no lost effect       & 12 \\ % operator.cfg
\bottomrule
\end{tabular}
\end{table}

\emph{Not checked, and the exclusions are the interesting half.} Every
cryptographic construction --- the request binding, the vault, the digest --- is
out of the model rather than modelled badly, because an uninterpreted function
standing in for a hash proves nothing about the real one and assuming it
unforgeable would assume exactly what is in question. The in-process
acknowledgement chain of the four links above is \emph{assumed}, not checked:
what the model checks is the part that is atomic and Redis-visible, the
authorization key and the preflight that requires it. Real time is not modelled;
where that cuts in the protocol's favour it is neutralised, since a lease may
expire at any instant, which is more adversarial than the budget, and where it
cuts against, it is the reconciliation delay named above. Liveness is checked
within the version bound, so it establishes that nothing gets stuck for a reason
internal to the protocol, not termination in an unbounded run.

\paragraph{Nine configurations exist to fail} Six configurations check that the
properties hold. The other nine each remove one stated assumption and must
produce the declared counterexample; a configuration that passes when it should
fail turns the build red, which is what keeps this from being a check that
cannot fail. \Cref{tab:modelchecking} is where the section's content actually
is. Two rows are worth reading twice. The AOF-rewind residual of
\cref{sec:p1} --- previously argued, then confirmed by a probe --- falls out in
four steps once the single-timeline assumption is dropped. And the
truthful-\texttt{fsync} row is the write-loss result of
\cref{sec:eval-durability} in the model: the barrier is issued, it reports
success, the write is discarded, and an effect ends up with no record at all.

\paragraph{The barrier is a query about durability, not a guarantee} The
ablated barrier of B3 returns success without asking anything. The real barrier
under write loss returns success because the server answers that it fsynced.
These are the same state in the model, and the specification expresses them with
one expression, which is the sharpest available statement of what the mechanism
is worth: \textbf{the barrier is not a durability guarantee but a question about
durability, and it inherits the truthfulness of whatever answers it}. The two
are not quite the same system --- the ablated barrier cannot report failure,
where the real one still fails when the server does not answer --- so write loss
removes the barrier's ability to inform while leaving its ability to fail.

\paragraph{What this establishes, and what it does not} TLC explored every
reachable state under these bounds and found no violation. That is a fact about
a specification of a few hundred lines, and four things bound what it transfers.
The model is \emph{not} the implementation: nothing connects them but a hand
transcription and one gate over the transition table and status alphabet, so
refinement is neither proven nor claimed, and a defect in the Lua's JSON
handling or arithmetic is invisible here. Redis is \emph{assumed}: script
atomicity, expiry and \texttt{WAITAOF} semantics are premises, and
\cref{sec:eval-durability} is the standing demonstration that one of them can be
false in a way nothing detects. The scope is \emph{small}: two workers, two
intents, five versions, one step, so a defect requiring a third worker or a
second step would not appear. And the properties not checked above are not
checked --- most of the implementation's security argument is among them. What
model checking changes is narrower than it looks and still worth having: the
composition of the fences, the state machine and the classification table was
previously an argument over code paths, and it is now a checked one --- together
with, in \cref{tab:modelchecking}, the price of each assumption it rests on.

\section{Implementation}
\label{sec:implementation}

AEP is 6\,849 lines of Python~3.13 against Redis~7.2.5, pinned by image digest.
% source: `find aep_core -name "*.py" | xargs wc -l`, 2026-08-07.
The evaluation harness, mock provider and six baseline systems are a further
20\,099 lines.
% source: `find experiments -name "*.py" -not -path "*/__pycache__/*" | xargs wc -l`.

\paragraph{What the protocol rests on} Three Lua scripts, each a single
serialisation point: the intent compare-and-swap that carries the whole
transition table, the pre-dispatch preflight, and the dispatch-authorization
writer. A Lua re-implementation of the strict UTF-8 and duplicate-JSON-member
checks is prefixed to every authoritative script, so the store rejects a
malformed document rather than trusting the client that sent it.

\paragraph{Composition modes} A start-up gate distinguishes \textsc{test},
\textsc{evaluation} and \textsc{production} compositions and refuses to run a
mismatched one. \textsc{production} and \textsc{evaluation} share one block of
requirements and differ only in the connector endpoint, which makes ``the
measurements ran in a production-shaped composition'' a property of the code
rather than of this sentence. \textsc{production} is nevertheless
\emph{unsatisfiable} in this repository --- there is no production vault and no
production connector, and the gate fails closed for it. Every number in
\cref{sec:evaluation} was therefore collected in \textsc{evaluation} mode, with
no test-only flags anywhere in the harness, and the evaluation vault's two
weaknesses relative to the production design (it shares the Redis trust domain,
and its keys are operator-supplied rather than KMS-wrapped) are declared rather
than hidden.

\paragraph{On test counts} The artifact carries a large automated suite and CI
that fails if any test is skipped. We deliberately do \emph{not} offer the test
count as evidence of protocol assurance. A prior internal audit of this
repository established that the suite had grown chiefly in rejection branches,
gate scripts and artifact metadata, while the number of tests exercising the
write-ahead path was close to unchanged; a reader who reads a doubled test
count as a doubled protocol guarantee would be wrong. The evidence for the
protocol's behaviour under faults is \cref{sec:evaluation}, and the harness
found three defects the suite had missed --- discussed in \cref{sec:threats}
because two of them would have biased results in AEP's favour had they gone
unnoticed.

\paragraph{Reproducibility} The environment is locked (\texttt{uv.lock}), the
Redis image is pinned by digest, and CI provisions Redis from the same compose
file the experiments use, so no job can pass against a differently configured
store. Continuous integration runs three jobs --- the full suite, the
\texttt{WAITAOF} durability suite, and a citation-range check over the formal
model --- and the build fails if any test is skipped.

\section{Evaluation}
\label{sec:evaluation}

\begin{itemize}
\item[\textbf{RQ1}] Under crashes, does AEP eliminate \emph{undetected}
      duplicates, and what does its residual ambiguity depend on?
\item[\textbf{RQ2}] Which of the protocol's two durability mechanisms produces
      which guarantee --- and against which faults is each one observable?
\item[\textbf{RQ3}] What does the protocol cost, and how much of that cost is
      a choice?
\item[\textbf{RQ4}] How does recovery behave?
\end{itemize}

\subsection{Setup}
\label{sec:eval-setup}

\paragraph{Systems} The seven of \cref{tab:trilemma}, implemented as thin
variants over one connector, one workload driver and one oracle, so that only
the differences differ. Each carries a machine-readable descriptor that tests
check against the implementation, and a crash point that does not exist in a
given system is marked inapplicable and never run rather than silently skipped.

\paragraph{Faults} Workers are separate OS processes killed with real
\texttt{SIGKILL} at one of the six crash points of \cref{tab:crashpoints}. A
\emph{regime} is a named fault condition --- not a matrix dimension --- and we
report each separately because they are different experiments:

\begin{itemize}
\item \textbf{crashed}: every execution is killed at the cell's crash point.
      This is the regime RQ1 is about.
\item \textbf{crash-free} (\texttt{p0}): nothing is injected. The only regime
      RQ3 may use.
\item \textbf{\texttt{redis-kill-preack}}: no worker crash; Redis is hard-killed
      with \texttt{docker kill -s KILL} at the instruction boundary between the
      intent CAS and the barrier acknowledgement, then restarted and verified.
      This is RQ2's ablation.
\end{itemize}

\paragraph{Oracle} A standalone mock provider records every \emph{actually
applied} mutation in a SQLite ledger (WAL, \texttt{synchronous=FULL}, one
transaction per mutation). Duplicates are counted from that ledger, never from
the system's own opinion of itself. The analysis is permitted to read exactly
two files per run --- the event log and the oracle ledger --- and a source gate
that parses the analysis modules fails the build if they import anything that
could reach live protocol state.

\paragraph{Statistics} Rates are over executions. Per-cell rate intervals use
a seeded run-cluster bootstrap (\BootstrapResamples{} resamples); the
B3--AEP ambiguity difference uses a stratified run-cluster bootstrap that
preserves crash-point, endpoint-capability, and keying strata. Fisher's exact
tests compare execution counts and therefore do not adjust for run clustering;
we report them descriptively, not as evidence of equivalence. Seeds and
repetition counts are recorded per run.

\paragraph{Two reporting rules we hold ourselves to} First, \emph{no pooled
table}. A rate computed across fault regimes is a property of how many runs of
each kind were collected, not of any system; our own analysis tool prints a
warning whenever more than one regime is present, and every number below names
its regime. Second, \emph{no absolute timing from an ungated host}. A run
contributes a duration only if the operator declared before the run that the
host cannot suspend \emph{and} the run's wall-versus-monotonic divergence stays
within tolerance. Both halves are necessary: a host that merely was not idle
long enough to suspend is indistinguishable from one that cannot. Under this
gate, \emph{zero} runs collected before this rule existed contribute a timing
number, and none does. Both rules were adopted after the first collection ---
the pooling rule when a mid-study review found the pooled summary table
unusable as a source, the timing gate after a host suspension was discovered
mid-run --- and every rate and duration below is reported under them.

\subsection{RQ1 --- undetected duplicates and the shape of the residual}
\label{sec:eval-rq1}

\Cref{tab:outcomes} is the anchor result; the supplementary material shows the
same executions pooled to one bar pair per system.

% GENERATED by scripts/paper_tables.py -- do not edit by hand.
% Source: analysis/per-cell-metrics.csv
% Filter: regime == '(session-3)' ONLY.
%   Every execution in these cells was killed at the cell's crash
%   point. The crash-free (p0) and hard-Redis-kill (redis-kill-preack)
%   regimes are DIFFERENT EXPERIMENTS and are reported separately;
%   pooling them is what got analysis/table-1.csv banned.
% Each cell = sum(successes)/sum(total) over the six crash points.
\begin{table*}[t]
\centering
\caption{The trilemma, measured. Every execution in every cell was killed at one of the six crash points of \cref{tab:crashpoints}; rates are over executions, pooled across crash points within one endpoint capability. \textsc{auth}/\textsc{pos-only}/\textsc{none} are the reconciliation capabilities of \cref{tab:capabilities}. AEP-full is the only system with a nonzero declared-ambiguity column, and the only one whose other two columns are zero everywhere. Source: \texttt{per-cell-metrics.csv}, crashed regime only.}
\label{tab:outcomes}
\small
\begin{tabular}{@{}lccccccccc@{}}
\toprule
& \multicolumn{3}{c}{undetected duplicate}& \multicolumn{3}{c}{lost effect}& \multicolumn{3}{c}{declared ambiguity}\\
\cmidrule(lr){2-4}\cmidrule(lr){5-7}\cmidrule(lr){8-10}
system & \textsc{auth} & \textsc{pos-only} & \textsc{none} & \textsc{auth} & \textsc{pos-only} & \textsc{none} & \textsc{auth} & \textsc{pos-only} & \textsc{none}\\
\midrule
B0 naive retry & 0.8200 & 0.8000 & 0.7733 & 0.0133 & 0.0133 & 0.0067 & 0.0000 & 0.0000 & 0.0000\\
B1 lease-only & 0.8133 & 0.8067 & 0.8267 & 0.0000 & 0.0067 & 0.0067 & 0.0000 & 0.0000 & 0.0000\\
B2 CAS-only & 0.8067 & 0.7933 & 0.7800 & 0.0067 & 0.0067 & 0.0000 & 0.0000 & 0.0000 & 0.0000\\
B4 durable, $\infty$ attempts & 0.9500 & 0.5889 & 0.5611 & 0.0000 & 0.0000 & 0.0111 & 0.0000 & 0.0000 & 0.0000\\
B4b durable, 1 attempt & --- & 0.0000 & 0.0000 & --- & 0.5056 & 0.5167 & --- & 0.0000 & 0.0000\\
B3 intent, no barrier & 0.0000 & 0.0000 & 0.0000 & 0.0000 & 0.0000 & 0.0000 & 0.0000 & 0.3667 & 0.7167\\
\textbf{AEP-full} & 0.0000 & 0.0000 & 0.0000 & 0.0000 & 0.0000 & 0.0000 & 0.0000 & 0.3500 & 0.7222\\
\bottomrule
\end{tabular}
\end{table*}

% Denominators (executions) per (system, response class):
%   B0_NAIVE_RETRY                     AUTHORITATIVE_READBACK   executions=150   runs=15   crash_points=5
%   B0_NAIVE_RETRY                     POSITIVE_ONLY_READBACK   executions=150   runs=15   crash_points=5
%   B0_NAIVE_RETRY                     NO_READBACK              executions=150   runs=15   crash_points=5
%   B1_LEASE_ONLY                      AUTHORITATIVE_READBACK   executions=150   runs=15   crash_points=5
%   B1_LEASE_ONLY                      POSITIVE_ONLY_READBACK   executions=150   runs=15   crash_points=5
%   B1_LEASE_ONLY                      NO_READBACK              executions=150   runs=15   crash_points=5
%   B2_CAS_ONLY                        AUTHORITATIVE_READBACK   executions=150   runs=15   crash_points=5
%   B2_CAS_ONLY                        POSITIVE_ONLY_READBACK   executions=150   runs=15   crash_points=5
%   B2_CAS_ONLY                        NO_READBACK              executions=150   runs=15   crash_points=5
%   B4_DURABLE_WORKFLOW                AUTHORITATIVE_READBACK   executions=20    runs=2    crash_points=1
%   B4_DURABLE_WORKFLOW                POSITIVE_ONLY_READBACK   executions=180   runs=18   crash_points=6
%   B4_DURABLE_WORKFLOW                NO_READBACK              executions=180   runs=18   crash_points=6
%   B4B_DURABLE_WORKFLOW_AT_MOST_ONCE  POSITIVE_ONLY_READBACK   executions=180   runs=18   crash_points=6
%   B4B_DURABLE_WORKFLOW_AT_MOST_ONCE  NO_READBACK              executions=180   runs=18   crash_points=6
%   B3_INTENT_NO_BARRIER               AUTHORITATIVE_READBACK   executions=180   runs=18   crash_points=6
%   B3_INTENT_NO_BARRIER               POSITIVE_ONLY_READBACK   executions=180   runs=18   crash_points=6
%   B3_INTENT_NO_BARRIER               NO_READBACK              executions=180   runs=18   crash_points=6
%   AEP_FULL                           AUTHORITATIVE_READBACK   executions=180   runs=18   crash_points=6
%   AEP_FULL                           POSITIVE_ONLY_READBACK   executions=180   runs=18   crash_points=6
%   AEP_FULL                           NO_READBACK              executions=180   runs=18   crash_points=6


Three readings, in the order they matter.

\paragraph{AEP's two silent columns are empty, everywhere measured} Across
\AepExecAuth{}, \AepExecPosOnly{} and \AepExecNoReadback{} crashed executions on
the three capability classes, AEP-full recorded an undetected-duplicate rate of
\AepDupAuth{}, \AepDupPosOnly{} and \AepDupNoReadback{}, and a lost-effect rate
of \AepLostAuth{}, \AepLostPosOnly{} and \AepLostNoReadback{}. Every one of
those executions was killed by \texttt{SIGKILL} at a named point around the
external call. The matrix runner halted on this condition: a single undetected
duplicate in AEP-full stops the matrix, a rule carried in the runner before any
of these cells were collected. We do not call it a pre-registration. The
governing document at the time described a contrary result as a defect to fix
rather than a finding to report, and this project's pre-registration practice
began later; what the runner gives is a stopping rule fixed in advance, not a
commitment about how the outcome would have been written up. It did not occur.

\paragraph{The baselines without a pre-dispatch record duplicate in most
crashed executions} B0, B1 and B2 sit between \BaselineDupLow{} and
\BaselineDupHigh{}, and the weakest of the three comparisons against AEP-full
is significant at $p = \BaselineDupMaxP{}$ by Fisher's exact test.
% \BaselineDupLow / \BaselineDupHigh: generated/numbers.tex, from
% per-cell-metrics.csv over {B0,B1,B2} x all collected response classes.
% p-value source: analysis/comparisons-vs-aep-full.csv, crashed regime.
The gap between B1/B2 and B0 is small and that is the finding: a lease and a
fenced CAS are doing exactly what they promise and neither promise is about the
external call (\cref{fig:bycrashpoint}).

\paragraph{A durable log is not enough, and the two configurations of one
engine land on two different corners} B4 has a write-ahead record,
acknowledged durable by the same \texttt{WAITAOF} barrier AEP uses. It is not
ablated on durability. It duplicates at \BfourDupPosOnly{} on
\textsc{pos-only} and \BfourDupNoReadback{} on \textsc{none}
(\BfourExecPosOnly{} and \BfourExecNoReadback{} executions), and it duplicates
on \textsc{auth} as well, at \BfourDupAuth{} over \BfourExecAuth{}
executions. The endpoint's
reconciliation capability does not rescue it, which is the point --- a
read-back tells a recovering workflow what the provider currently holds, and
B4's duplicate is a second application it chose to make, not a fact it failed
to look up.

B4b is the same engine with Maximum Attempts set to $1$, and it is the cleaner
half of the demonstration: over \BfourbExecPosOnly{} and
\BfourbExecNoReadback{} executions it records \emph{zero} undetected
duplicates on both classes and loses the effect in \BfourbLostPosOnly{} and
\BfourbLostNoReadback{} of them instead, and \BfourbLostAuth{} over
\BfourbExecAuth{} on \textsc{auth}. One configuration line moves
the engine from one silent corner of \cref{sec:trilemma} to the other, and
neither configuration reaches the third: across all \BfourFamilyAmbCells{}
B4 and B4b cells the declared-ambiguity rate never exceeds
\BfourFamilyAmbMax{}.

\paragraph{The engine itself was run, and it separates what B4 licenses from
what it does not} We ran the production engine against the same non-cooperative
endpoint, in the two configurations B4 and B4b instantiate, at the one crash
point both the engine and our model can be cut at. With unlimited attempts it
duplicated at \BfiveDup{} over \BfiveExec{} executions and lost effects at
\BfiveLost{}; with Maximum Attempts set to one it lost at \BfivebLost{} over
\BfivebExec{} executions and duplicated at \BfivebDup{}. Neither configuration
declared ambiguity in any of \BfiveRuns{} non-void runs (\BfiveAmb{} and
\BfivebAmb{} declarations) --- the third corner of \cref{sec:trilemma} going
unreached by the engine itself and not merely by our model of it.

\emph{The direction transfers; the magnitudes do not.} Each configuration lands
in the corner the supplementary comparison assigns it, with exactly zero events
of the other kind. But at that same crash point B4 duplicates at
\BfourAtBarrier{} over \BfourAtBarrierExec{} executions and B4b loses at
\BfourbAtBarrier{} over \BfourbAtBarrierExec{}, roughly six times the engine's
rates. \textbf{B4's rates are our model's, not the product's}, and no rate in
this paper should be read as an engine's.

Four things bound this and none is removed by having run the engine. The
comparison is \emph{incomplete}: the engine has no position answering to
\textsc{after-intent-before-barrier}, because its pre-dispatch record is written
by the server inside one RPC, so there is no worker-side instant between the
write and its acknowledgement --- a hole exactly where B4 and B4b have data at
all three capability classes. The rates are conditional on a
\BfiveStartToClose{}\,ms Start-To-Close timeout that B4 has no analogue for; a
different timeout is a different cell, and no timing number here is any engine's
recovery latency. The interval arithmetic is asymmetric --- B4's intervals rest
on three runs against the engine's thirty, and three clusters put resampled
means on multiples of one thirtieth --- so the separation should be read from
the point estimates, which differ by roughly $0.8$, and not from the intervals
failing to overlap. And the crash window at the dispatch point is sized by the
server and by loopback rather than by the protocol, so its width is not
like-for-like with B4's.

That is the sharpest form of the argument the paper can make \emph{against our
model of a durable-execution engine} --- and that qualification is load-bearing
rather than modest, and the paragraphs above are why: B4 is our implementation
of the pattern, not a product, and we registered before collecting B5 that a
gap of that size would oblige us to re-scope every B4 claim. The shape of the
argument transfers and the magnitude is our model's. The argument is nonetheless
made against a system that has everything the write-ahead pattern is supposed to
provide. The record is necessary and it is not sufficient, because the
record cannot contain the fact that matters --- whether the provider applied
the effect. What separates B3 and AEP-full from B4 is not durability but the
\emph{absence of an edge back into the dispatching state}: the transition
table has no path that re-enters \textsc{about-to-fire}, checked inside the
atomic script (\cref{fig:statemachine}). What separates them from B4b is that
declining to re-send is not the same as knowing whether to.

\subsubsection{The residual is set by the endpoint, not the protocol}

AEP's declared-ambiguity column is \AepAmbAuth{} on \textsc{auth},
\AepAmbPosOnly{} on \textsc{pos-only} and \AepAmbNoReadback{} on
\textsc{none}. That ordering is not a defect that better engineering would
remove; it is \cref{tab:capabilities} showing through. Where the endpoint can
prove absence, recovery resolves everything and nothing is left ambiguous.
Where it can only confirm presence, absence is uninformative and the honest
answer is ``I do not know''. Where it cannot be queried at all, no amount of
protocol recovers the fact.

\paragraph{The \textsc{pos-only} rate depends on how the read-back is keyed, and
we did not expect it to} The read-back in the cells above is keyed by the
caller's own reference. We pre-registered an alternative keying --- an oracle
fingerprint of the request --- as a sensitivity check, and predicted it would
leave every headline rate inside a $\pm\BthreeVsAepAmbMargin{}$ percentage-point
margin. On \textsc{pos-only} it does not: declared ambiguity is
\KeyingAmbiguityCaller{}\% under the caller's reference and
\KeyingAmbiguityOracle{}\% under the fingerprint, a difference of
\KeyingAmbiguityDelta{} percentage points, and the margin was
\BthreeVsAepAmbMargin{}.
% generated/numbers.tex | ws5-2026-09-10/t2-keying/analysis/
% per-cell-metrics.csv against analysis/per-cell-metrics.csv, metric=
% known_ambiguity_rate, system=AEP_FULL, response_class=POSITIVE_ONLY_READBACK,
% pooled over the six crash points as sum(successes)/sum(total).

The two columns that carry the safety claim are unmoved: undetected duplicates
and lost effects are zero under both keyings. What moves is the residual, and
the consequence is a limit on generality rather than on correctness. **The
\textsc{pos-only} ambiguity rate is a property of the protocol \emph{and} of how
the recovering worker phrases its question**, so the figure quoted above should
be read as this keying's, not as the class's. \textsc{auth} is unaffected,
because a read-back that can prove absence resolves the execution whichever key
it is asked under.

The supplementary material plots the same executions as one bar pair per system,
pooled over crash points and capabilities: the three systems on the right have an
empty left bar, and only AEP-full and its ablation B3 have a right bar at all.
The figure adds no number this section does not state; the bars are Wilson
intervals on pooled counts, and the cluster-aware bootstrap intervals in
\texttt{per-cell-metrics.csv} are wider and are the ones to quote.

The rate decomposes further by crash point, and that is where the protocol's
conservatism becomes visible; the supplementary material gives the full
decomposition across all six crash points and three capability classes.

Two cells carry the argument. At \texttt{before\_intent\_write} the rate is
zero in every capability class: the crash precedes the intent, nothing was
promised, and there is nothing to be ambiguous about. At
\texttt{after\_barrier\_before\_dispatch} it is at its maximum --- and that is
the crash point at which \emph{no effect can possibly exist}, because the
worker died before sending anything. AEP declares ambiguity most often exactly
where it is least warranted by the world. That is the false-positive ambiguity
of \cref{sec:p2}, and its direction is the whole design: the protocol's errors
are toward saying ``I might have done something'' when it did not, never toward
silence when it did. The cells where an effect \emph{does} typically exist ---
\texttt{mid\_dispatch} and \texttt{after\_response\_before\_resolution} under
\textsc{pos-only} --- show low ambiguity, because a positive read-back
confirms the effect and the intent resolves.

\subsubsection{Why the provably-empty cell is not resolved}
\label{sec:eval-provable}

The \texttt{after\_barrier\_before\_dispatch} row invites an obvious question:
if the worker died before sending anything, \emph{the system could know that}.
Why is the outcome ambiguous rather than a confirmed non-event?

Because nothing durable records it. A recovery process reading the store
afterwards sees an authorized, unresolved intent --- exactly what it would see
had the worker died \emph{during} transmission, with the request already on the
wire. The two histories are identical in everything durable, so recovery fails
closed, which is the correct behaviour given what it can observe. \textbf{That
is the negative result, and it is the protocol's, not the implementation's:} the
cell is empty in the world and unresolvable from the store. Closing it is
possible and we did not do it. The supplementary material gives the probe detail
and prices the third barrier that would split the two cases --- a trade that
narrows rather than closes the window, and whose value depends entirely on a
deployment's ratio of crash points, which our uniform matrix deliberately does
not model.

% Single column rather than full width. A float's cost is its area, and for a
% graphic -- unlike a table -- narrowing genuinely shrinks the area rather than
% re-wrapping the same rows taller: measured at 0.29 pages saved, against 0.02
% to 0.04 for the three tables that were tested the same way and left alone.
\begin{figure}[t]
\centering
\includegraphics[width=\columnwidth]{figures/figure-2-duplicates-by-crash-point.pdf}
\caption{Undetected duplicate rate by crash point, crashed regime, pooled over
endpoint capabilities weighted by executions. The baselines' rate is not
concentrated at one crash point: any interruption after the request is on the
wire is enough, which is why a mechanism that acts only \emph{after} the call
cannot help.}
% Generated by experiments/analyze.py write_figures(), regime "(session-3)".
\label{fig:bycrashpoint}
\end{figure}

\subsubsection{The cost of the third corner}

Declared ambiguity is not free, and a paper that presented it as a pure win
would be selling something. Each declared ambiguity is an execution that has
stopped, that will not proceed without a human, and whose resolution requires
someone to establish out-of-band what the endpoint will not say. On a
\textsc{no-readback} endpoint under a crash-heavy regime that is most of the
crashed executions. What we claim is not that this is cheap but that it is
\emph{the honest price of the information the endpoint withholds}, and that
the alternative is not fewer incidents but the same incidents discovered later
by a different department. We have not evaluated whether operators handle
declared ambiguities well, or how much a real deployment would generate; both
are stated as gaps in \cref{sec:threats}.

\subsection{RQ2 --- which mechanism produces which guarantee}
\label{sec:eval-rq2}

AEP contains two durability mechanisms and the literature usually treats them
as one. The first is the \emph{pre-dispatch record}: an intent, written and
fenced, that exists before the external call does. The second is the
\emph{barrier}: the \texttt{WAITAOF} round trip that makes the record's
durability a checked precondition of dispatch authority (\cref{sec:p2}). B3 is
the ablation that separates them --- AEP-full with the barrier removed and
nothing else removed --- and running it against three different fault classes
assigns each guarantee to the mechanism that actually produces it.

The assignment is not the one we expected, and it is the paper's central
structural result.

\subsubsection{Detection is the record's, not the barrier's}
\label{sec:eval-detection}

Everything in \cref{sec:eval-rq1} --- zero undetected duplicates, zero lost
effects, a residual of declared ambiguity shaped by endpoint capability ---
survives the barrier's removal intact.

% GENERATED by scripts/paper_tables.py -- do not edit.
% Source: analysis/per-cell-metrics.csv (rates, regime=crashed)
%         analysis/comparisons-vs-aep-full.csv (regime-labelled p-values).
\begin{table}[t]
\centering
\caption{The barrier ablation on the detection metrics. B3 is AEP-full with the \texttt{WAITAOF} barrier removed and nothing else removed. Crashed regime; rates are over executions, pooled across crash points within one capability class. The p-values are execution-level Fisher tests over all capability classes; they do not account for run clustering and are not used as equivalence evidence.}
\label{tab:ablation}
\small
\begin{tabular}{@{}llccc@{}}
\toprule
metric & capability & AEP-full & B3 & $p$\\
\midrule
\multirow{3}{*}{undetected duplicate} & \textsc{auth} & 0.0000 & 0.0000 & \multirow{3}{*}{$1.00$}\\
 & \textsc{pos-only} & 0.0000 & 0.0000 & \\
 & \textsc{none} & 0.0000 & 0.0000 & \\
\midrule
\multirow{3}{*}{lost effect} & \textsc{auth} & 0.0000 & 0.0000 & \multirow{3}{*}{$1.00$}\\
 & \textsc{pos-only} & 0.0000 & 0.0000 & \\
 & \textsc{none} & 0.0000 & 0.0000 & \\
\midrule
\multirow{3}{*}{declared ambiguity} & \textsc{auth} & 0.0000 & 0.0000 & \multirow{3}{*}{$0.95$}\\
 & \textsc{pos-only} & 0.3500 & 0.3667 & \\
 & \textsc{none} & 0.7222 & 0.7167 & \\
\bottomrule
\end{tabular}
\end{table}


Over \BthreeVsAepN{} crashed executions per arm, B3 and AEP-full show no
observed difference in either silent-failure metric. Both record an
undetected-duplicate count of $\BthreeVsAepDupCount{}$ and a lost-effect count
of $\BthreeVsAepLostCount{}$. Declared ambiguity is
\BthreeAmbCount{}/\BthreeVsAepN{} for B3 and
\AepAmbCount{}/\BthreeVsAepN{} for AEP-full; per capability class the two
systems track each other: \BthreeAmbAuth{} against \AepAmbAuth{} on \textsc{auth},
\BthreeAmbPosOnly{} against \AepAmbPosOnly{} on \textsc{pos-only},
\BthreeAmbNoReadback{} against \AepAmbNoReadback{} on \textsc{none}.

For ambiguity, B3 minus AEP-full is \BthreeVsAepAmbDiffPP{} percentage points.
The stratified run-cluster \BthreeVsAepAmbDiffConfidence{}\% interval is
[\BthreeVsAepAmbDiffLow{},\,\BthreeVsAepAmbDiffHigh{}] percentage points,
which lies within a $\pm\BthreeVsAepAmbMargin{}$ percentage-point margin. We
introduce that margin in this revision; it was not preregistered, and we state
it as a stipulation rather than derive it. We do not claim it is immaterial to
an operator: \cref{sec:threats} records that we have not evaluated declared
ambiguity as an operational outcome, and that the artifact has no escalation
mechanism, so we have no basis on which to price one. The result
supports equivalence under that stated, post hoc margin, but only for this
fixed design. Each arm has \BthreeVsAepAmbClusters{} run clusters distributed
over \BthreeVsAepAmbStrata{} matched strata---three runs per stratum---so the
cluster-aware interval is a sensitivity analysis, not a strong general
equivalence result. The execution-level Fisher value
($p = \BthreeVsAepAmbP{}$) is not used as equivalence evidence.

\textbf{A word on what the zeros can and cannot support.} Fisher's exact test
on $\BthreeVsAepDupCount{}/\BthreeVsAepN{}$ against
$\BthreeVsAepDupCount{}/\BthreeVsAepN{}$ returns
$p = \BthreeVsAepDupP{}$ for duplicates and $p = \BthreeVsAepLostP{}$ for lost
effects. Neither $p$-value carries information: a test between two zero
counts has no power, and failing to distinguish them is not evidence that they
are the same. What two zeros do support is a \emph{bound}. The unit for that bound
is the run, not the execution: each arm's \BthreeVsAepN{} executions are
\BthreeVsAepRuns{} runs of ten, and ten executions that share a run, a crash
point and a system are not ten independent trials. The one-sided
95\% Wilson upper confidence bound on a zero numerator over
\BthreeVsAepRuns{} run clusters is \AblationZeroUpperRun{}\%. This is a genuinely
one-sided bound using the 95th normal quantile; it is not the upper endpoint of
a two-sided 95\% Wilson interval, which uses the 97.5th quantile. Each system's
rate is individually bounded this way. When the two system bounds are stated
simultaneously, Bonferroni gives joint coverage of at least 90\%, not 95\%.
Under that simultaneous event, both rates lie in
$[0,\AblationZeroUpperRun{}\%]$, so their absolute difference is at most
\AblationZeroUpperRun{} percentage points. This finite-sample bound is the support
for ``no observed difference'' in the two zero-event metrics; the corresponding
Fisher values are not evidence of equivalence.

Treating the executions as independent would give
\AblationZeroUpperExec{}\%. We use the run, for the same reason the Fisher
values in \cref{tab:ablation} are not read as equivalence: with no events
observed, nothing in these data estimates the within-run correlation that
execution-level independence would require. The choice of unit decides nothing
claimed below:
the baseline duplicate rates this bound is contrasted against are
\BaselineDupLowPct{}--\BaselineDupHighPct{}\%, more than an order of magnitude
from either figure, so the conservative unit is taken because it is the
defensible one and not because the comparison needs it.

\textbf{We report this as a finding rather than as a limitation, because it
reassigns our own headline claim.} The property that distinguishes AEP from
B0, B1 and B2 --- and the property the paper is named for --- is produced by
the pre-dispatch record and by the transition table that refuses to
re-enter \textsc{about-to-fire}. It is not produced by waiting for
\texttt{fsync}. A reader who wants the detection guarantee and not the
barrier's cost can have it, and \cref{sec:eval-deployment} prices that
configuration explicitly.

This also sharpens what B4 shows. B4 has a durable record, acknowledged by the
same barrier, and duplicates at the highest rate in the study; B3 has a durable
record, no barrier, and duplicates at zero. Durability is neither necessary nor
sufficient for detection. What is necessary is that the record be written
\emph{before} the call and that no path re-enter the dispatching state.

\subsubsection{Prevention is the barrier's, and it needs the right fault}
\label{sec:eval-prevention}

If detection does not need the barrier, what does the barrier do? It withholds
the external call when the durability of the intent can no longer be confirmed.
That is a different guarantee, it is worth stating in its own metric, and no
worker-crash fault can exercise it --- which is why \cref{sec:eval-detection}
returns a null. The fault that exercises it is loss of Redis itself.

The metric is the \emph{unwanted-applied-effect rate}: the fraction of
executions in which a real, non-idempotent mutation reached the provider while
the durability of its own intent record could no longer be established. It is
not the same quantity as an undetected duplicate and it is not measured on the
same regime; keeping the two apart is the whole point of this subsection.

We hard-kill Redis at \texttt{after\_intent\_before\_barrier}: the instruction
boundary between the intent CAS and the acknowledgement. AEP-full is inside
\texttt{WAITAOF} at that moment; B3 is not, because B3's ablation is precisely
that round trip. One execution per run, no worker crash --- the fault under
study is Redis dying.

\paragraph{The controlled fault} The result below is collected under a
\emph{controlled} injector, and that is the change that makes it worth
reporting. \texttt{docker kill} alone lands \UncontrolledLanding{}\,ms after it
is issued on this host; freezing the container first with \texttt{docker pause}
and killing it afterwards lands in \ControlledLanding{}\,ms, with a third of the
spread, and leaves the container in the same end state. The fault class does not
change --- it is the same F3 crash-stop at the same instruction boundary --- so
the two are comparable, and the injector's own contribution to the race is
narrowed by measurement rather than argued away.

Over \ArmASessions{} sessions and all \ArmAClasses{} capability classes,
AEP-full's unwanted-applied rate is \ArmAAepRate{} (\ArmAAepApplied{} of
\ArmAAepN{} executions) and B3's is \ArmABthreeRate{}
(\ArmABthreeApplied{}/\ArmABthreeN{}), on the same injected fault. The
between-session spread in AEP-full's count --- the quantity a single collection
cannot see, and the one the earlier version of this result was criticised for
not reporting --- collapses to \ArmASpreadMax{} count per class across the
\ArmASessions{} sessions. B3 stays at its ceiling throughout, never below
\ArmABthreeFloor{} of \AepKillRuns{} in any class in any session, which is what
rules out the reading that the injector simply disabled both arms.

\paragraph{The result fell below its own pre-registered bound, and we report
that as a surprise rather than a success} The prediction was registered before
the data existed: \ArmAPredicted{}, bounded [\ArmABoundLow{},
\ArmABoundHigh{}], derived from the measured landing distribution as the
fraction of a $1\,000$\,ms \texttt{WAITAOF} window the injector leaves open.
The observed \ArmAAepRate{} is \emph{below} that bound. The pre-registration
fixed in advance what a rate outside the bound would mean --- the protocol
behaving unexpectedly, not the mechanism succeeding --- and we honour that here
rather than presenting a lower-than-predicted number as a better result.

The reason the bound was wrong is measurable, and it is not calibration: the
model behind it fails on two independent premises that push in opposite
directions. First, it treats the exposure window as the injector's bench-measured
landing latency. That does not survive contact with a collection:
\ModelAcksPastLanding{} of \ModelAckN{} uncontrolled AEP-full runs received a
genuine \texttt{WAITAOF} acknowledgement \emph{after} the bench landing had
supposedly passed, and a \texttt{WAITAOF} acknowledgement requires a live
server. Timed during the collections themselves, the same injector call takes
\ModelInSituKill{}\,ms against the bench's \ModelBenchKill{}\,ms ---
\ModelInSituRatio{}$\times$, with ranges that do not overlap. Second, it treats
the \texttt{WAITAOF} wait as uniform on $U(0, 1\,000)$\,ms. Measured over the
runs the model was calibrated against, the fraction of waits under
\ControlledLanding{}\,ms is \ModelEmpiricalWindow{}, where uniformity assigns
\ModelUniformWindow{} --- \ModelUniformOverstates{}$\times$ too generous, and
worse the narrower the window. The two errors partly cancelled at the
uncontrolled window and could not at the controlled one.

The calibration that appeared to license the bound compared the wrong two
quantities. The model predicts \emph{dispatch}; the estimand is an
\emph{applied effect}, and dispatch is necessary for one without being
sufficient. On the same runs the dispatch rate is \ModelDispatchRate{} and the
applied rate is \ModelAppliedRate{}. Against the quantity it actually predicts,
the model was already wrong on its own calibration data. We record this because
a pre-registered bound that agrees for the wrong reason is worse than one that
disagrees, and only the second kind is visible without looking.

\paragraph{What this does and does not license us to say} We report a rate and
its mechanism, not a guarantee. \textbf{We do not claim that AEP-full prevents
the unwanted effect.} The residual is non-zero and is irreducible under this
fault class: freezing Redis synchronously at the checkpoint would make
\texttt{WAITAOF} unanswerable, but B3 reaches the same checkpoint and still
needs two further Redis calls, so the same freeze would stop B3 dispatching too
and the contrast would vanish because the injector had disabled both arms. What
the controlled fault buys is that the residual is now bounded by something we
measured rather than by this host's \texttt{docker} latency --- and, as above,
the bound we derived from that measurement was itself wrong in a way the data
can show. The magnitude is not a constant of the protocol and this paper does
not offer it as one.

\paragraph{The uncontrolled cell replicates the direction, and is not pooled
with the controlled one} The original collection used \texttt{docker kill}
alone, and is kept as the uncontrolled replication. It is reported separately
rather than pooled: the two injectors have different landing distributions, and
combining them would trade a stated limitation for an unstated confound. There
the rate is \AepUnwantedRate{} for AEP-full against \BthreeUnwantedRate{} for B3
(Fisher's exact test, two-tailed, $p = \UnwantedP{}$), a difference of
\UnwantedPrevented{} executions on the same injected fault --- the same
direction as the controlled cell, at a magnitude the uncontrolled injector sets.
This is the only fault in the study under which the two systems diverge;
everywhere else they differ by at most $\BthreeVsAepAmbDelta{}$ executions in
\BthreeVsAepN{} (\cref{tab:ablation}). Both systems declared ambiguity on all
\AepKillRuns{} runs: the endpoint is \textsc{no-readback}, so neither can
resolve what it does not know, and both fail closed. The difference between them
is not in what they \emph{say} --- that is the detection guarantee, and
\cref{sec:eval-detection} already showed it does not depend on the barrier ---
it is in what they did to the outside world before saying it.

\textbf{The direction replicates; the magnitude moves, which is why the
controlled cell exists.} The cell was re-collected \ReplicationSessions{}
further times under an identical configuration, and AEP-full's unwanted-applied
count ranged \ReplicationAepMin{}--\ReplicationAepMax{} out of \AepKillRuns{}
while B3 stayed at \ReplicationBthreeApplied{}/\AepKillRuns{} in every one.
Taking the session as the unit, the difference averaged
\ReplicationPreventedMean{} effects, 95\% interval
[\ReplicationPreventedLow{},\,\ReplicationPreventedHigh{}]; against that spread
the controlled cell's \ArmASpreadMax{}-count per-class spread over
\ArmASessions{} sessions is the change worth reporting. Two negative findings
travel with the cell and are not softened by being moved: attributing the
movement to kill latency did not succeed at this design's precision, and B3 is
not the negative control an earlier version of this paper reported it to be ---
its own session-level interval is far too wide to have contradicted the
mechanism whatever it had done. The per-session tables carrying both --- the
uncontrolled counts, the four session-level kill-latency differences with their
intervals, and the same analysis for B3 --- are in the supplementary material.
The scope of the cell is one pre-acknowledgement crash point and one host, and
the \textsc{authoritative-readback} class was subsequently collected over
\ClassSessions{} sessions (\cref{sec:threats}) at a precision inadequate to the
question.

\paragraph{Where the fault lands after the branch point, the arms tie} The
same fault delivered \emph{after} dispatch separates nothing, and we collected
that rather than inferring it. Arming the kill at the last instruction before
transmission and delivering it $200$\,ms later puts it in flight, after both
arms have already dispatched. Over \InflightSessions{} sessions the two arms
differ by at most \InflightMaxDiff{} count in any capability class, and both sit
at their ceiling. Two qualifications belong with that number. The second session
is a deterministic replay of the first rather than an independent draw --- the
harness assigns seeds per cell and repetition, so where no race remains the
second collection reproduces the first exactly, which is itself the clearest
evidence that no race remains. And B3's lowest cell is
\InflightBthreeFloor{}/\AepKillRuns{}, exactly the floor below which the result
would have been read as a mis-timed injector rather than a tie; it is the
boundary, not headroom.

What is structural, across all three faults, is that the race exists for
AEP-full at all and cannot exist for B3.

\subsubsection{The barrier's durability claim, under the fault that can
actually test it}
\label{sec:eval-durability}

There remains a third thing the barrier is supposed to do, and it is the one
most readers assume it does first: protect the write-ahead record from being
lost when Redis dies. We tested that premise before building on it,
and for the fault everyone reaches for first, it is false.

\ProcessKillTrials{} trials, each phase-aligning the \texttt{everysec} timer
with a \texttt{WAITAOF} so the window under test is as wide as it can be,
writing a key \emph{without} waiting for it, hard-killing the container with
\texttt{docker kill -s KILL}, restarting, and reading back. Every kill landed
(every container came back with its \texttt{uptime} at zero) and every kill
landed inside the window (write-to-death
\ProcessKillWindowMin{}--\ProcessKillWindowMax{}\,ms against a $1\,000$\,ms
fsync period).
\textbf{Not one unfsynced write was lost: $\ProcessKillUnackLost{}$.}
% source: reports/raw/e1-durability-window.txt; Session 3B section C.2.

The mechanism is not subtle once seen. \texttt{appendfsync everysec} defers the
\texttt{fsync(2)}, \textbf{not} the \texttt{write(2)}: Redis writes the AOF
buffer to the operating system on every event-loop iteration. A
\texttt{SIGKILL} destroys the \emph{process}; the kernel and its page cache are
untouched, and a kernel that is still running flushes those bytes to disk. The
same result reappears at $n{=}\KillCanaryN{}$ inside the ablation cells above, where a
canary write made immediately before each kill survived in \AepKillCanary{} and
\BthreeKillCanary{} runs.

That narrows the claim to a named fault class --- loss of the page cache: host
power failure, kernel panic, VM destruction --- and an earlier draft of this
paper stopped there, naming a fault it had not injected. We now inject it.

\paragraph{Making the storage lose the writes} The probe is the one above with
the fault swapped. A loop-backed block device carries a \texttt{dm-flakey}
target whose \texttt{drop\_writes} feature silently discards every write bio
while serving reads correctly, which is what a storage stack looks like to a
kernel whose writes will never reach the platter. Redis 7.2.5 --- the same
build as the rest of the evaluation --- runs with its append-only file on an
ext4 filesystem over that device. Each trial writes one key through
\texttt{WAITAOF} and one without, switches the device into
\texttt{drop\_writes}, kills the server, unmounts to discard the page cache,
restores the device and reads both keys back. The switch is the instant of
power loss: what was \texttt{fsync}-ed before it is on the device, what was
still dirty is not.

Two details decide whether this measures anything. \texttt{dmsetup suspend}
freezes the filesystem and flushes outstanding I/O by default, either of which
would push the unacknowledged write out before the fault took effect and
produce a confident null; the reload runs with \texttt{--nolockfs --noflush}.
And a self-test runs first on the device alone, with no Redis in the picture: a
file \texttt{fsync}-ed before the switch must survive and one \texttt{fsync}-ed
after it must not. The probe refuses to report Redis trials unless that holds.
% source: experiments/flakey_write_loss.py; reports/raw/g2-flakey-write-loss.txt.

\paragraph{The result} The probe ran \FlakeyN{} trials in \FlakeyReplications{}
independent replications. Each replication builds the device stack from
scratch, and each \emph{trial} within it rebuilds the filesystem with
\texttt{mkfs.ext4} and starts a fresh server, so no trial can read a key a
previous trial's append-only file still holds. \textbf{Separation was perfect in
every replication}: the acknowledged record survived
\FlakeyPerRepAckSurvived{} and the unacknowledged one was destroyed
\FlakeyPerRepUnackLost{} in each, \FlakeyAckSurvived{} and \FlakeyUnackLost{}
overall. Against the same two keys under a process kill,
$\ProcessKillUnackLost{}$ were lost; the fault classes separate completely.

We give no $p$ for either comparison, and the reason is that no unit supports
one: at the trial level the test would assume an independence the shared device
stack violates and would ignore the pairing --- both records are written in the
same trial and meet the same write-loss event --- while at the replication
level, three against three, the two-sided minimum a perfectly separated result
can reach is $0.1$, which is the design's floor rather than a statement about
these data. The complete separation in every replication is the finding, and it
does not need one. The filesystem was mounted
\texttt{ext4 rw,relatime} with its journal enabled and no \texttt{nobarrier},
and the server's \texttt{appendfsync} was read back as \texttt{everysec} from
the running instance; the probe refuses to measure if either is otherwise,
because under \texttt{always} there is no unacknowledged write to lose and
under \texttt{nobarrier} a surviving acknowledged one would prove nothing.

\textbf{The barrier's durability benefit to the \emph{record} is therefore real
and now measured, and the fault class matters more than the mechanism.} What
that benefit is worth to a \emph{dispatch} is the next subsubsection's question,
and the answer there is not the one we expected. The Redis documentation's
``you may lose 1 second of data if there is a disaster''~\cite{redis-persistence}
is correct; what our two probes add is that ``disaster'' must mean the host,
that the widely repeated reading of the sentence as covering a crashed Redis
\emph{process} is wrong, and that once the host is the fault the sentence is
exactly right.

One qualification, and it is the same one the process-kill probe carries: the
exposure window is deliberately maximised. The unacknowledged write is issued
\FlakeyWindowMin{}--\FlakeyWindowMax{}\,ms before the device stops accepting
writes, inside a $1\,000$\,ms fsync period that a preceding \texttt{WAITAOF}
has just reset. \FlakeyUnackLost{} is therefore the loss probability for a
power cut placed where it does the most damage, not the rate a uniformly-timed
one would produce; that rate is bounded above by it and approaches it as the
write approaches the start of the period. The design is deliberate and it is
the same design that made the process-kill probe's $\ProcessKillUnackLost{}$
strong: both probes
give the loss hypothesis its best case, and only one of them finds loss.

\subsubsection{Running the protocol under that fault, and what it refuted}
\label{sec:eval-writeloss-cell}

The probe above tests \texttt{WAITAOF}, not AEP. It runs two Redis keys against
a device that stops accepting writes; it does not run AEP-full and B3 through
the harness under that fault, and so it measures no difference in \emph{protocol}
outcomes. What it establishes is the premise the durability argument rests on.
We then ran the extension, and it is reported here rather than as a caveat
because it refuted what we predicted.

\paragraph{What was registered, and what happened} We pre-registered that
AEP-full would withhold dispatch when its write-ahead record was destroyed and
that B3 would proceed. A full session was collected with the harness's Redis
running its append-only file on the flakey device: both arms,
\WriteLossRunsPerArm{} runs each, \WriteLossExecPerRun{} executions per run,
\texttt{drop\_writes} armed at the intent checkpoint, without a delivery
failure. \textbf{The prediction is refuted.} AEP-full applied an effect in
\WriteLossAepApplied{} of \WriteLossExecPerArm{} executions and B3 in
\WriteLossBthreeApplied{}; both arms finished at ceiling, and neither withheld.

\paragraph{Why, and it is the more useful finding} \textbf{\texttt{WAITAOF}
returned success for every durability acknowledgement requested after the device
had stopped accepting writes} --- \WriteLossAcksAfterFault{} of them, with
\WriteLossAckFailures{} failures. \texttt{dm-flakey}'s \texttt{drop\_writes}
discards each write bio and reports success, so the kernel sees no error,
\texttt{fsync(2)} succeeds, and Redis has no way to learn otherwise. The barrier
withholds dispatch on a \emph{failed or absent} acknowledgement; it was handed a
successful one and dispatched, behaving exactly as specified on an input that
was false.

\paragraph{The bounds, stated narrowly} \textbf{This is not evidence that the
barrier fails to withhold under write loss}: the condition the experiment exists
to create --- the barrier meeting an acknowledgement that fails or never arrives
--- was never reached, so no run exercised withholding. \textbf{Neither is it
evidence that the barrier works.} What it establishes is that the durability
guarantee rests on an unstated premise --- that a \texttt{WAITAOF} success means
the record reached the device --- and that under storage which discards writes
\emph{silently} the premise is false and the protocol cannot detect it. Storage
returning \texttt{EIO} would let the acknowledgement fail and put the barrier
back on the input its guarantee is defined over, so this result bounds the
guarantee to storage that reports its own failures and says nothing about
storage that does not. It is evidence about storage, not about the protocol. One
host, one Redis version (7.2.5, \texttt{appendfsync everysec}), and we claim it
for no other.

\subsection{RQ3 --- cost, and how much of it is optional}
\label{sec:eval-rq3}

Only crash-free runs on a suspend-disabled host contribute. The provider is
configured with a constant $2\,000$\,ms delay, so one round trip is the floor
and the interesting quantity is the increment over it.

% GENERATED by scripts/paper_tables.py -- do not edit.
% Source: analysis/latency-and-throughput.csv
% 'median' is step_latency_ms_median over crash-free runs only
% (denominator = overhead_runs_crash_free, printed).
% 'over B0' subtracts B0's median (2010.165771 ms), which is
% the no-protocol system on the same host and the same provider
% (configured delay 2000 ms; B0 sits above it by
% 10.2 ms, which is harness cost, not protocol cost).
% E5 gate: a run contributes a duration only if the host was declared
% non-suspending BEFORE the run and no wall-versus-monotonic
% divergence was observed in it. Counts are unaffected.
\begin{table}[t]
\centering
\caption{Median step latency, crash-free runs only, E5-gated. Increments are over B0, the no-protocol system on the same host and provider. The whole write-ahead protocol minus the barrier is the B3 row; the barrier is the difference between B3 and AEP-full.}
\label{tab:latency}
\small
\begin{tabular}{@{}lrrr@{}}
\toprule
system & runs & median step (ms) & over B0 (ms)\\
\midrule
B0 no protocol & 3 & 2\,010.2 & ---\\
B1 $+$ lease & 3 & 2\,014.0 & 3.9\\
B2 $+$ fenced CAS & 3 & 2\,017.4 & 7.3\\
B3 full protocol, no barrier & 3 & 2\,038.2 & 28.0\\
\textbf{AEP-full} & 3 & 4\,004.9 & 1\,994.7\\
B4b durable, 1 attempt & 3 & 4\,013.4 & 2\,003.3\\
B4 durable, $\infty$ attempts & 3 & 4\,015.7 & 2\,005.5\\
\bottomrule
\end{tabular}
\end{table}


\textbf{The overhead is the barrier, and the rest is not separable from
zero.} A median step costs \AepStepMedian{}\,ms under AEP-full against
\BzeroStepMedian{}\,ms under B0, which does nothing but call the provider.
Almost all of that gap is one mechanism. At fifteen crash-free runs per arm the
two \texttt{WAITAOF} round trips cost \BarrierCostFifteen{}\,ms between them,
with a 95\% interval of
[\BarrierCostFifteenLow{}, \BarrierCostFifteenHigh{}]\,ms.
% generated/numbers.tex | ws5-2026-09-10/t1-p0-everysec/analysis/
% per-execution.csv, cluster bootstrap over runs, 10000 resamples, seed
% 20260806: median AEP-full - median B3.

Everything the write-ahead protocol does apart from waiting for \texttt{fsync}
--- the intent ledger, the vault, the request binding, the preflight, the
recovery service --- costs
\ProtocolMinusBarrierFifteen{}\,ms
[\ProtocolMinusBarrierFifteenLow{}, \ProtocolMinusBarrierFifteenHigh{}]\,ms on
a two-second call. That interval spans zero. B3's arm is a mixture of two
modes, and restricting both arms to their lower mode gives
\ProtocolMinusBarrierLowerMode{}\,ms
[\ProtocolMinusBarrierLowerModeLow{}, \ProtocolMinusBarrierLowerModeHigh{}]\,ms,
which spans zero as well. We report both because they were both pre-registered
and they disagree, and choosing between them after seeing them is the move this
paper is arguing against.
% generated/numbers.tex | same CSV and bootstrap: median B3 - median B0,
% pooled over both modes, and again restricted to each arm's lower mode. The
% lower-mode boundary is the splitter's own lower_max, fixed before this arm
% was collected.

\textbf{So the decomposition supports an ordering and not a partition.} The
barrier dominates: its interval is nowhere near zero and the residual's
contains it under both readings. What cannot be said is what fraction of the
protocol's cost is optional, because at this precision the non-barrier
remainder is not distinguishable from nothing at all. An earlier draft of this
paper quoted that remainder as a point figure of about thirty milliseconds; it
rested on three runs per arm, and twelve more did not confirm it.

Read against \cref{sec:eval-detection} the ordering still says something
stronger than ``AEP costs two seconds''. Whatever the detection guarantee
costs, it is the difference between B0 and B3 --- B3 is the system that already
has that guarantee in full --- and it is small enough that fifteen runs cannot
separate it from zero. The \BarrierCostFifteen{}\,ms buys the prevention
guarantee of \cref{sec:eval-prevention} and nothing else.

That figure is also a property of a \emph{configuration}, not of AEP. Under
\texttt{appendfsync everysec} an acknowledgement cannot return until the next
scheduled fsync, and those are one second apart, so a barrier costs up to a
second and the protocol pays that twice. The measured total sits near the top
of that range rather than in the middle of it; the reason is unmeasured.

\subsubsection{The barrier is a deployment choice, with three measured points}
\label{sec:eval-deployment}

The $\approx 2$-second figure above is the one a reader would carry away as
``AEP's overhead''. It is not: it is the price of one of the protocol's two
guarantees, under one of the durability settings that guarantee can be bought
at. \Cref{tab:deployment} is the decision as a deployment actually faces it.

% GENERATED by scripts/paper_tables.py -- do not edit.
% Sources: analysis/latency-and-throughput.csv (everysec rows),
%          fsync-always/analysis/latency-and-throughput.csv (always).
% step_latency_ms_median over crash-free, E5-gated runs only.
% Barrier column = row median - the B3 median collected under the
% SAME fsync policy (everysec 2\,038.2 ms, always 2\,048.4 ms),
% i.e. the same protocol with the barrier ablated and nothing else
% changed. Never across policies: that would assume the ablated
% protocol's own writes cost the same under both, which is exactly
% what the two B3 numbers are here to establish rather than assume.
\begin{table*}[t]
\centering
\caption{The barrier is a deployment choice, not a fixed cost. All three rows run the same pre-dispatch intent ledger and therefore make the same detection claim; they differ in what they pay to also withhold dispatch when durability cannot be confirmed. The barrier column is each row's own median minus a B3 median collected under the same \texttt{appendfsync} policy. Crash-free runs only, E5-gated, 2\,000\,ms provider floor.}
\label{tab:deployment}
\small
\renewcommand{\arraystretch}{1.08}
\begin{tabularx}{\textwidth}{@{}>{\raggedright\arraybackslash}Xrrrc>{\raggedright\arraybackslash}p{0.20\textwidth}@{}}
\toprule
configuration & \shortstack{median\\(ms)} & \shortstack{over\\floor} & \shortstack{barrier\\(ms)} & prevents & claim\\
\midrule
AEP-full, \texttt{everysec} & 4\,004.9 & 2\,004.9 & 1\,966.7 & yes & detection + prevention\\
AEP-full, \texttt{always} & 2\,063.4 & 63.4 & 15.0 & yes & detection + prevention\\
B3-mode, \texttt{everysec} & 2\,038.2 & 38.2 & 0.0 & no & detection only\\
\midrule
\multicolumn{6}{@{}p{0.96\textwidth}@{}}{\footnotesize The detection claim of \cref{tab:outcomes} shows no observed difference down the whole table: it is produced by the pre-dispatch record plus no re-entry, which all three rows have, and \cref{tab:ablation} is the ablation that shows it. What the barrier buys is the last column's second word, and \cref{tab:killablation} is what it is worth. `Over floor' is the same median less the provider's 2\,000\,ms delay, and so includes the 28.0\,ms the protocol costs with the barrier already removed.}\\
\bottomrule
\end{tabularx}
\end{table*}


Three points, and the third is the one the earlier framing of this paper
concealed.

\paragraph{\texttt{everysec} with the barrier} The default, and the expensive
one. An acknowledgement must wait for the next scheduled fsync, up to a second
away, and the protocol waits twice.

\paragraph{\texttt{always} with the barrier} The identical crash-free cells ---
same systems, endpoint, regime, seeds, host and provider --- against a Redis
whose only changed line is \texttt{appendfsync always}.
% source: reports/raw/f0iii-appendfsync-always.txt and the G2 re-run that
% added the B3 arm; both refuse to proceed unless CONFIG GET appendfsync
% reads back 'always' on the live server.
The barrier's cost is \BarrierCostAlwaysFortyFive{}\,ms
[\BarrierCostAlwaysFortyFiveLow{}, \BarrierCostAlwaysFortyFiveHigh{}]\,ms over
fifteen runs per arm: negative, and not distinguishable from zero. Under
\texttt{always} the fsync has already happened by the time the write returns,
so \texttt{WAITAOF} has nothing left to wait for.
% generated/numbers.tex | fsync-always-2026-09-14/analysis/
% latency-and-throughput.csv and per-execution.csv, same bootstrap as above.

\textbf{That is the measurement this section turns on.} The barrier does not
become cheap under \texttt{always}; it costs nothing, and the
\BarrierCostFifteen{}\,ms it costs under \texttt{everysec} is therefore the
price of the fsync boundary rather than of the barrier. A protocol step under
\texttt{everysec} waits for the next scheduled flush, and there are two of
them. Nothing in the mechanism is expensive.

The sign is not a puzzle to explain away. The two medians it separates are
\AepAlwaysFortyFiveMedian{}\,ms with the barrier and
\BthreeAlwaysFortyFiveMedian{}\,ms without it, on a two-second call; the gap
between them is noise around a mechanism that is doing no waiting. This arm was
pre-registered with the interval's exclusion of zero explicitly ruled out as an
acceptance criterion, because a mechanism with nothing to wait for should cost
nothing and a measurement of nothing lands on either side of it.
% generated/numbers.tex | fsync-always-2026-09-14/analysis/
% latency-and-throughput.csv, step_latency_ms_median per system.

Both figures are ablation differences \emph{within} one policy, which required
collecting the barrier-less arm twice. It would have been easier to subtract
the \texttt{everysec} B3 median from the \texttt{always} AEP-full median, and
that would have assumed the ablated protocol's own writes cost the same under
both policies. They nearly do --- B3 is \BthreeStepMedian{}\,ms under
\texttt{everysec} and \BthreeAlwaysFortyFiveMedian{}\,ms under
\texttt{always} --- but
``nearly'' is a measurement, and an order-of-magnitude claim resting on an
unmeasured premise is the kind of number this paper is trying not to publish.

\paragraph{How precise these two figures are} Each is a difference of two
medians over fifteen runs per arm, bootstrapped over runs rather than
executions. Under \texttt{everysec} the barrier costs
\BarrierCostFifteen{}\,ms, [\BarrierCostFifteenLow{}, \BarrierCostFifteenHigh{}];
under \texttt{always} it is \BarrierCostAlwaysFortyFive{}\,ms,
[\BarrierCostAlwaysFortyFiveLow{}, \BarrierCostAlwaysFortyFiveHigh{}]. The first
excludes zero by a wide margin and the second contains it, which is the whole
of the claim: the policy decides whether the barrier is expensive, and the
barrier itself is not.

\paragraph{B3-mode: the barrier removed entirely} Not a baseline we are beating
but a supported configuration of the same implementation, licensed by
\cref{sec:eval-detection}. What it pays over a system with no protocol at all
is not resolvable from zero at fifteen runs per arm
(\ProtocolMinusBarrierFifteen{}\,ms,
[\ProtocolMinusBarrierFifteenLow{}, \ProtocolMinusBarrierFifteenHigh{}]). On the detection metrics, pooled across capability
classes, we observe no difference between it and AEP-full, supported by a bound
rather than a test: both the undetected-duplicate and lost-effect rates lie in
$[0,\AblationZeroUpperRun{}\%]$ at joint coverage of at least 90\%, and the
declared-ambiguity difference lies inside the $\pm\BthreeVsAepAmbMargin{}$
percentage-point margin \cref{sec:eval-detection} states as a post hoc
stipulation. \emph{We do not make that claim per capability class, and the
reason is a number rather than an oversight}: six simultaneous bounds instead of
two, over a third as many run clusters, puts the Bonferroni width at
\AblationZeroUpperPerClass{} percentage points, and a band that wide is not a
detection guarantee. What B3-mode forfeits is \cref{sec:eval-prevention}: under
a hard Redis kill it will put an effect on the wire that AEP-full withholds, at
\BthreeUnwantedRate{} against \AepUnwantedRate{}, and its intent record will not
survive a host-level write loss (\cref{sec:eval-durability}). That second
clause is about the \emph{record}, and it does not license the protocol-level
inference a reader may draw from it: when we ran both arms under exactly that
fault, neither withheld, because the storage discarded the writes while
reporting success and the barrier was handed an acknowledgement that looked
good (\cref{sec:eval-writeloss-cell}).

\textbf{So the question a deployment answers is not ``can we afford AEP'' but
``do we need prevention, and at which fsync policy''.} If losing Redis is the
fault to worry about --- an instance on the same host as the workers, an
unreplicated instance, a VM that can vanish --- the barrier is the point of the
protocol and \texttt{always} makes it nearly free at this workload's shape. If
Redis is well protected and worker crashes are the concern, B3-mode gives the
paper's headline result for a cost this evaluation cannot distinguish from zero.

\textbf{One measurement this table does not support} is what \texttt{always}
costs under load. The Redis
documentation calls it ``very very slow, very safe''~\cite{redis-persistence}, a
judgement about write throughput that our workload cannot exercise: two workers,
a two-second provider delay and a handful of writes per step leave the system
latency-bound and nowhere near an fsync-rate limit. Our throughput figure is in
fact \emph{higher} under \texttt{always} ($\AepThroughputAlways{}$ versus
$\AepThroughputEverysec{}$ executions/s), which is a statement about this
workload's shape and not about Redis.

\subsection{RQ4 --- recovery}
\label{sec:eval-rq4}

Recovery behaviour is reported through the outcome distribution rather than a
success rate, for a reason worth stating: ``recovery success'' means two
different things either side of whether a recovery service was running. Where
one ran, a crashed execution reaching a terminal classification was
\emph{recovered}. Where none ran --- B0, B1, B2, B4, B4b declare none --- the
same execution reached its classification because a supervisor ran the step
again, which is not recovery but the thing recovery exists to avoid. Our
analysis excludes runs without a recovery service from that metric rather than
scoring them zero or one, because both would be statements about a component
that was not present.

Within that scope: every AEP-full and B3 execution in the crashed regime
reached a terminal classification, and across \RunsCollected{} runs exactly
one recorded a reconciliation disagreement --- every other completed run's
event log agreed with its independently written oracle ledger. The harness
performs this reconciliation after each completed run; it checks definitive
applied/non-applied assertions, effect-count bounds, attribution, and duplicate
predictions rather than trusting a line count.
% source: analysis/coverage.json and the per-run reconcile step.

That one is worth stating precisely, because the check exists to be believed.
It was a B4b cell in which the system classified all ten executions
\textsc{confirmed-applied} while the provider's ledger recorded two, eight
disagreements of the form ``the system asserted the effect was applied and the
ledger records none''. Both sibling runs of the same cell recorded ten, the
provider logged no error, and its run-log held three lines against eleven for
each sibling --- so the provider recorded almost nothing while answering the
workers successfully. We could not determine the cause. A re-run of the
identical cell, same seed and same configuration, agreed.

We treat it as a \emph{void} run on the rule the durability probe uses: a
trial whose measurement precondition failed says nothing about the system in
either direction. Counting it would let a broken ground truth set a baseline's
rate; discarding it quietly and keeping the re-run would be re-running until
the number looks right. So the re-run is what \cref{tab:outcomes} contains,
the voided attempt must ship in the external raw archive under
\texttt{results/voided/} with this explanation beside it; it is not present in
the current Git repository. The count above is one rather than zero.
Recovery latency is not reported as an absolute figure: the E5 gate leaves too
few gated crashed runs to support a distribution, and \cref{sec:threats}
records this as a gap rather than filling it with ungated numbers.

\section{Related Work}
\label{sec:related}

AEP composes established mechanisms: leases, fenced compare-and-swap,
write-ahead logging, and an fsync barrier. We do not claim that durable intent
plus locking is new. The contribution is a capability-specific recovery
semantics for effects outside the coordinator's atomicity domain, and an
evaluation that separates detection from prevention.

\paragraph{Leases and fencing} Leases provide time-bounded mutual exclusion
under crash faults~\cite{gray1989leases}; Chubby made the pattern an
infrastructural service~\cite{burrows2006chubby}, and the Redis lock
pattern~\cite{redis-locks} is a commodity form. A lease holder may stall past
expiry and write afterwards, so a monotonic fencing token is also
required~\cite{kleppmann2016locking}. AEP applies that standard remedy to its
state store. B1 and B2 show its boundary: fencing protects the store, not an
external endpoint that does not interpret the token (\cref{tab:outcomes}).

\paragraph{Logging, outboxes, and long-lived transactions} The
expected-version CAS is optimistic concurrency control~\cite{kung1981optimistic}.
ARIES writes intent before an action that the same recovery domain can undo or
redo~\cite{mohan1992aries}; AEP instead records an action performed by a
foreign endpoint that may support neither operation. The transactional-outbox
pattern likewise stores a message before a relay publishes it, but explicitly
permits the relay to publish twice after a crash and therefore assumes an
idempotent consumer~\cite{richardson-transactional-outbox}. B3 is structurally
outbox-like, with a no-re-entry rule in the measured runner; it does not make
the remote mutation transactional. Sagas are appropriate when a compensating
action exists~\cite{garciamolina1987sagas}. AEP addresses the case in which no
compensation exists, or deciding whether to compensate first requires the
missing outcome.

The expired IETF HTTPAPI working-group draft for an \texttt{Idempotency-Key}
header documents the cooperative alternative: a server can make non-idempotent
\texttt{POST} and \texttt{PATCH} requests safely retryable by interpreting a
client key~\cite{jena2025idempotencykey}. AEP's premise is precisely the
endpoint that does not implement such a contract. Helland describes the
broader consequences of operating without distributed transactions and the
centrality of idempotence~\cite{helland2007apostate,helland2012idempotence}.

\paragraph{Locks with intent and durable execution} Olive is the closest
mechanistic antecedent. Its locks with intent durably associate code with a
lock and provide exactly-once execution despite client failures and concurrent
duplicate executions~\cite{setty2016olive}. A write-ahead intent as a checked
precondition of execution is therefore prior work, not AEP's novelty. Olive's
guarantee is realized through operations made recoverable within its cloud
storage abstraction. AEP starts where that atomicity scope ends: the effect is
performed by a legacy endpoint, and in the \textsc{no-readback} class no
participant can later establish whether it happened. AEP consequently offers
no exactly-once claim; it terminates in declared ambiguity and prohibits
automatic re-dispatch.

Beldi~\cite{zhang2020beldi}, Durable
Functions~\cite{burckhardt2021durable}, Boki~\cite{jia2021boki}, and
ExoFlow~\cite{zhuang2023exoflow} develop transactional or replay-based durable
execution. Their guarantees are appropriate when an effect is idempotent,
transactional, compensable, or otherwise recoverable. Temporal states the
operational premise directly: activities should be idempotent because they may
execute more than once~\cite{temporal-activity-definition}; setting Maximum
Attempts to one selects the opposite retry policy~\cite{temporal-retry}, and
the engine's own account of when an activity is considered failed rests on
timeouts and heartbeats rather than on any fact about the
endpoint~\cite{temporal-activity-failures}. Our B4 and B4b controls instantiate
these two policies against the same non-cooperative endpoint. They are semantic
controls, not claims about those products' complete implementations.

That distinction is no longer only a caveat. We subsequently ran the production
engine itself against the same endpoint, in both configurations, at the one
crash point the engine and our controls can both be cut at
(\cref{sec:threats}). \emph{The direction transferred and the magnitudes did
not}: unlimited attempts duplicated and lost nothing, one attempt lost and
duplicated nothing --- each landing in the corner \cref{tab:related} assigns it
--- while the rates were several times below B4's and B4b's at the same crash
point, and neither configuration declared ambiguity in any run. The comparison
covers one crash point of five, is conditional on a fixed activity timeout the
controls have no analogue for, and rests on frozen control cells an order of
magnitude thinner than the engine's; it is reported with those bounds in
\cref{sec:threats} and no rate in this paper should be read as a product's.

The ablation supplies AEP's narrower empirical delta. B3 and AEP-full share
the pre-dispatch record and no-re-entry transition but differ on the durability
barrier. We observe no difference in their crashed-regime detection outcomes,
bounded by pooling the capability classes: both zero-event rates fall within
$[0,\AblationZeroUpperRun{}\%]$ at joint coverage of at least 90\%
(\cref{sec:eval-detection}). The Redis-kill experiment isolates a prevention
effect. We found no prior evaluation that measures this detection/prevention
decomposition with an independent external-effect ledger.

\paragraph{Agent execution reliability} Three contemporaneous systems overlap
substantially and delimit the claim. Verified Tool Calls formalizes non-atomic
agent tool failures and uses postcondition verification, verify-before-retry,
and idempotency keys in a simulated fault environment~\cite{mansoor2026verified}.
Its verifier requires a queryable postcondition; without server-side
idempotency, it characterizes the remaining race as best-effort. AEP evaluates
the authoritative, positive-only, and absent-readback cases separately and
uses terminal ambiguity for the last case.

LogAct durably exposes intentions in a shared log before execution and supports
recovery and policy voting for open-ended agent control flow
~\cite{balakrishnan2026logact}. It is broader in agent-state-machine structure
and isolation; AEP is narrower in studying what can be concluded about one
non-idempotent external mutation after an interruption. The Sovereign Execution
Broker (SEB) reserves certificate nonces before mutation, disallows retry after
possible invocation unless durable or target-side evidence proves otherwise,
and provides a mandatory broker boundary with cryptographic certificates
~\cite{he2026seb}. That boundary is stronger than AEP's trusted-process guard.
AEP's distinct evidence is the three endpoint-capability classes, terminal
declared ambiguity, real \texttt{SIGKILL} and Redis durability faults, and the
B3/AEP ablation. ACRFence addresses semantic rollback during agent
checkpoint-restore~\cite{zheng2026acrfence}; it is adjacent rather than a basis
for ranking which work is ``nearest.''

\paragraph{Durability acknowledgements and what they are worth} AEP's barrier
withholds dispatch until the store acknowledges the intent durable, so the
protocol is only as strong as that acknowledgement. A body of storage work
establishes that the acknowledgement can be false. Rebello et al.\ show that
none of five widely deployed applications --- PostgreSQL, LMDB, LevelDB, SQLite
and Redis --- handles \texttt{fsync} failure sufficiently, and that on Linux the
affected pages are marked clean regardless of whether writeback
succeeded~\cite{rebello2020fsync}; the operational discovery that a successful
\texttt{fsync} does not mean the data reached persistent storage was made in
PostgreSQL and reported at the time~\cite{corbet2018fsync}. Ganesan et al.\
find that eight distributed storage systems, Redis among them, do not
consistently use redundancy to recover from file-system faults, so a single
fault can produce data loss or unavailability~\cite{ganesan2017redundancy}.
Alagappan et al.\ show that recovery which ignores protocol-level knowledge
loses data in replicated logs, and repair it for
Raft-style~\cite{ongaro2014raft} systems~\cite{alagappan2018par}.

We place our own durability result in that line rather than beside it:
\cref{sec:threats} reports a barrier that dispatched on an acknowledgement which
was already impossible to honour, because the underlying device discarded the
write and reported success. The protocol behaved exactly as specified on an
input that was false. What this literature supplies is the reason that is a
property of the storage stack rather than a defect peculiar to our
implementation --- and the reason a barrier's guarantee must be stated as
conditional on the acknowledgement it rests on.

\paragraph{Exactly-once state versus exactly-once effect} Consistent recovery
for distributed computation is classically obtained by snapshotting global
state~\cite{chandy1985snapshots}, and modern dataflow engines make that
practical by injecting barriers into the stream and persisting operator state
asynchronously~\cite{carbone2015abs}. The guarantee these mechanisms provide is
that \emph{internal state} is restored to a consistent cut. An external
mutation already performed is not inside that cut and is not undone by
restoring it, which is why exactly-once processing is routinely paired with a
transactional or idempotent sink. AEP's endpoint is neither, and in the
\textsc{no-readback} class it cannot even be interrogated after the fact --- so
the snapshot line of work and AEP are solving adjacent problems on opposite
sides of the atomicity boundary. This is also the end-to-end
argument~\cite{saltzer1984endtoend} in an uncomfortable form: the check that
would settle whether the effect happened belongs at the endpoints, and the
endpoint that would have to perform it is precisely the one that cannot.

\paragraph{Fault-injection lineage} AEP's named crash boundaries and
independent oracle follow a mature crash-consistency methodology. ALICE
constructs possible post-crash storage states and checks application
invariants~\cite{pillai2014alice}; Torturing Databases combines controlled
power-fault injection with record/replay and cross-layer diagnosis
~\cite{zheng2014torturing}; and CrashMonkey/Ace exhaustively test bounded
workloads at persistence boundaries~\cite{mohan2018crashmonkey}. AEP applies
this lineage to a caller-to-endpoint protocol: real process kills at named
instruction boundaries, a separately written provider ledger, and a
\texttt{dm-flakey} write-loss probe. The harness is an evaluation instrument,
not a claimed fault-injection contribution.

\begin{table*}[t]
\centering
\caption{Selected design assumptions. ``Endpoint must cooperate'' denotes
idempotence, an idempotency key, transaction enrolment, or compensation. The
rows are not ordered by guarantee strength; they expose different residuals
when their stated assumptions fail.}
\label{tab:related}
\small
\renewcommand{\arraystretch}{1.08}
\begin{tabularx}{\textwidth}{@{}>{\raggedright\arraybackslash}X>{\raggedright\arraybackslash}p{0.19\textwidth}c>{\raggedright\arraybackslash}p{0.18\textwidth}@{}}
\toprule
 & endpoint must & can report & residual on \\
 & cooperate & \emph{unknown} & failure \\
\midrule
Lease / fenced CAS~\cite{gray1989leases,kleppmann2016locking}
  & n/a & no & duplicate effect \\
Sagas~\cite{garciamolina1987sagas}
  & compensation & no & compensation gap \\
Transactional outbox~\cite{richardson-transactional-outbox}
  & idempotent consumer & no & duplicate delivery \\
\texttt{Idempotency-Key} header~\cite{jena2025idempotencykey}
  & key support & no & duplicate effect \\
Locks with intent~\cite{setty2016olive}
  & effect inside storage & no & --- \\
Dataflow snapshots~\cite{chandy1985snapshots,carbone2015abs}
  & transactional sink & no & duplicate external effect \\
Beldi~\cite{zhang2020beldi}
  & yes & no & --- \\
Durable Functions~\cite{burckhardt2021durable}
  & yes & no & duplicate or loss \\
Boki~\cite{jia2021boki}
  & yes & no & --- \\
ExoFlow~\cite{zhuang2023exoflow}
  & yes & no & --- \\
Temporal-style, $\infty$ retries~\cite{temporal-retry}
  & yes & no & duplicate effect \\
Temporal-style, 1 attempt~\cite{temporal-retry}
  & yes & no & lost effect \\
\midrule
\textbf{AEP} & \textbf{no} & \textbf{yes} & \textbf{declared ambiguity} \\
\bottomrule
\end{tabularx}
\end{table*}

\paragraph{Revised novelty position} AEP's claim is not the invention of
write-ahead intent, locking, or fail-closed interposition. It is the measured
composition for endpoints with authoritative, positive-only, or no readback:
capability-specific recovery; a terminal declared-ambiguity outcome when no
query can settle the effect; real process and durability-fault evaluation with
an independent ground-truth ledger; and an ablation that assigns detection to
the pre-dispatch record plus no re-entry, and prevention to the durability
barrier.

% ORDER. Artifact Availability sits BEFORE Threats to Validity so the paper's
% last numbered section ends on an argument rather than on a URL. §VIII's
% closing subsection, "Where the protocol should not be used", states the
% precondition under which a durable-execution engine is the better tool and
% is the last thing a reader is told.
%
% The file names keep their original numbers -- 08-threats renders as §IX and
% 09-artifact as §VIII. Renaming the files would churn every path in the
% reports without changing a rendered word, and nothing reads the numbers:
% every cross-reference in the manuscript is a \cref on a label, and the
% supplementary names main-paper sections in words because it cannot \cref
% across documents (check_paper_numbers.py::check_cross_document_references).
\section{Artifact Availability}
\label{sec:artifact}

The public repository \artifactavail{} contains the implementation, evaluation
harness, mock provider, baseline systems, analysis pipeline, manuscript source,
and the tracked derived analysis products used by the number generator. It does
\emph{not} contain the raw run directories, which are bulk data: those are
carried by a separate archive, \archiveavail{}.

\paragraph{The raw evidence archive} Two archives under one DOI: the
2026-09-03 build (20 collection roots, 1\,458 run directories, 26\,300
files) and a 2026-09-15 extension covering the trees it did not reach (11
roots, 1\,332 run directories, 18\,494 files) --- 2\,790 run directories and
44\,794 files in all, 848\,MB uncompressed and 42\,MB compressed. Between
them they carry the 432-run evaluation, the \texttt{appendfsync=always} arm,
the deployment sweep of \cref{sec:eval-rq3}, the real-Temporal baseline, the
write-loss cell, the voided oracle-disagreement run, every prevention
replication, and the arms voided for having run against the wrong container
runtime. Every file carries a SHA-256 in its own archive's
checksum manifest, and that manifest's own digest ---
\texttt{87fa2d53\ldots} and \texttt{54d1ab0f\ldots} --- attests that archive;
the two share no path and each is verified only against its own. Both were
extracted and checked file-by-file before deposit, 26\,300 and 18\,494
entries, none failing or missing. The 2026-09-03 archive was additionally
re-analysed --- 114 tracked analysis products byte-identical, 8 identical
after two declared normalisations recorded with the deposit, none differing
--- and \emph{that re-analysis has not been repeated for the extension}.
% Numbers discipline (main.tex header). These are properties of the deposit,
% not of a tracked CSV, so they are not macros: their sources
% (ARCHIVE-METADATA.json and MANIFEST.sha256 inside each archive) live
% outside the repository, and a generated macro reading them would fail in
% the clean clone an evaluator uses. Each figure below names its source.
%   20 roots / 1 458 run dirs / 26 300 files / 87fa2d53... / 3aa90b21... :
%     aep-raw-archive/{ARCHIVE-METADATA.json,MANIFEST.sha256},
%     docs/29 1, ARTIFACT.md 5, reports/phase-report-11-rescue-2026-09-03.md
%   11 roots / 1 332 run dirs / 18 494 files / 54d1ab0f... :
%     aep-raw-archive-ext/{ARCHIVE-METADATA.json,MANIFEST.sha256},
%     docs/29 0b, ARTIFACT.md 5,
%     reports/phase-report-35-archive-extend-2026-09-15.md
%   848 MB = 492 905 568 + 355 492 262 payload bytes; 42 MB = 24 257 505 +
%     17 563 367 compressed. 2 790 = 1 458 + 1 332; 44 794 = 26 300 + 18 494.
%   A run directory is one holding a single run's artifacts. The extension's
%     ARCHIVE-METADATA.json reports 992, not 1 332, because its counter
%     counts run-config.json files and 340 directories in four roots -- the
%     three B5/Temporal session attempts and one voided write-loss arm --
%     carry run artifacts without one, B5 running through session.py rather
%     than run_matrix. docs/29 3 discloses this to a reader who opens the
%     metadata. Counted by scripts/verify_published_archive.py.
%   114 / 8 / 0: scripts/verify_raw_archive.py over the 2026-09-03 archive
%     only, most recently in the phase-39 clean clone. The extension has no
%     equivalent run, which is why the sentence says so.

\paragraph{What is pinned} The Python environment is locked
by a resolved dependency lockfile; the Redis image is pinned by digest in the
compose file;
the matrix records a seed per run and emits its full plan, cost estimate and
platform fingerprint \emph{before} launching. Continuous integration
provisions Redis from the same compose file the experiments use, runs the full
suite, the \texttt{WAITAOF} durability suite, a citation-range check over
the formal model and the numbers gate described below, and fails if any test is
skipped.

\paragraph{What can be reproduced, and at what cost} The matrix is resumable
and its cells are independent. The full plan is 1\,068 runs at an estimated
25\,h of wall time on one machine; the cells behind the results in
\cref{sec:evaluation} are a subset keyed in the tracked per-cell manifest. The raw archive carries the complete per-cell and
per-run manifest alongside the runs, so a reader can see what was collected
rather than inferring it from the tables.

\paragraph{Regenerating the paper's numbers} The tables and every headline
scalar in \cref{sec:evaluation} are generated by
the artifact's table generator from the tracked inputs its documentation
enumerates: per-cell and per-execution metrics, latency under both
fsync policies, the regime-labelled comparisons, the Redis-kill ablation,
coverage, and the write-loss probes. Each emitted value carries a comment
naming the filter and arithmetic behind it.
The artifact's build script builds the manuscript and runs that
check; it treats an undefined reference and a BibTeX parse error as build
failures rather than warnings, because both have produced a clean-looking
build of this paper once each.
An independent number check re-derives the tables and fails if
the manuscript has drifted from the results. It also fails if a generated
number is defined and never used, which is the drift a framing revision
produces:
LaTeX catches a macro used without a definition and is silent about a claim
whose evidence was orphaned when the claim moved. The pooled summary table
produced by the analysis tool is deliberately \emph{not} a source for any
number here. All comparisons used by the paper are filtered to an explicit
regime; the headline and B3 comparisons use only the crashed regime.

\paragraph{The host-level write-loss probe} \Cref{sec:eval-durability} needs a
block device that stops accepting writes, which needs root, a loop device and
a \texttt{dm-flakey} target, so the write-loss probe is
not part of the test suite --- a test that skips wherever those are
unavailable is the failure mode this project's CI gate exists to prevent. What
\emph{is} in the suite is everything about the probe that can be checked
without a kernel: that its ``drop'' table really carries \texttt{drop\_writes},
that its pass-through table carries no feature at all, and that a trial whose
acknowledged write vanished is voided rather than counted. The probe itself
refuses to report any Redis trial until a self-test has shown, on the bare
device, that a file \texttt{fsync}-ed before the switch survives and one
\texttt{fsync}-ed after it does not.

\paragraph{The state-machine figure} \Cref{fig:statemachine} is generated by a script
that imports the transition set from the implementation. An edge added to or removed from the protocol changes the
figure or fails the check; it cannot silently disagree with the code.

\paragraph{Citations} Every bibliography entry was verified to exist before it
was written: matched against DBLP's search API, or resolved through
\texttt{doi.org}, or --- for vendor documentation --- fetched and read on a
recorded date. The verification script and bibliography metadata are tracked.
The raw lookup transcript is not present in the Git repository and is not in
the raw-evidence archive either, which carries run data only; it is the one
item named in this section for which no external copy yet exists.

\section{Threats to Validity}
\label{sec:threats}

% Organised as construct / internal / external / conclusion validity, one threat
% stated once. Material enumerated rather than argued -- the coverage gaps of the
% fault-injection matrix, and the measurement incidents found during collection
% -- is in the supplementary material, pointed to from the subsection it belongs
% to. check_paper_numbers.py scans paper/supplementary.tex as well as
% main.tex and sections/*.tex, so a macro that moved with that material is still
% checked for use.

\subsection{Construct validity}

\paragraph{The barrier's durability benefit is measured on an emulated fault,
not a real power cut} \Cref{sec:eval-durability} injects the fault class
\texttt{WAITAOF} defends against --- loss of the page cache --- with a
\texttt{dm-flakey} target in \texttt{drop\_writes} mode, and the result is
unambiguous ($\FlakeyAckSurvived{}$ acknowledged records survive,
$\FlakeyUnackLost{}$ unacknowledged ones do not). What remains is that
\texttt{drop\_writes} is an \emph{emulation}. It is faithful in the dimension
that matters --- writes issued after a chosen instant never reach the platter,
while everything \texttt{fsync}-ed before it does --- and unfaithful in others:
a real power loss can tear a sector mid-write, can lose writes the device
claimed to have committed to its own volatile cache, and does not stop at a
boundary we chose. Those differences make real hardware \emph{worse} than the
emulation, so the direction of the result is safe and its precision is not: we
cannot claim the acknowledged record would survive a real power cut on a device
with a lying write cache, because \texttt{fsync} there does not mean what the
barrier assumes. That is a property of the storage, which AEP inherits from
every system that trusts \texttt{fsync}. The probe also maximises the exposure
window by phase-aligning to a just-completed fsync, so $\FlakeyUnackLost{}$ is
the loss rate for a worst-placed cut rather than a uniformly-timed one.

\paragraph{The write-loss probe tests \texttt{WAITAOF}, not AEP}
\Cref{sec:eval-durability} runs two Redis keys against a device that stops
accepting writes. It does not run AEP-full and B3 through the evaluation harness
under that fault, and so it does not measure a difference in \emph{protocol}
outcomes --- duplicates, lost effects, declared ambiguity --- attributable to
write loss. What it establishes is the premise the durability argument rests on:
that an acknowledged record survives an event that destroys an unacknowledged
one, which under a process kill is false and under write loss is true.

\paragraph{We ran the extension, and the fault is deliverable while the
durability signal above it is not trustworthy} The harness's Redis now runs with
its append-only file on the flakey device, and a full session --- both arms,
thirty runs each, ten executions per run, \texttt{drop\_writes} armed at the
intent checkpoint --- was collected without a delivery failure, so
\textbf{the fault is deliverable} for a whole session on this host. The second
half is answered in a way we did not predict. \textbf{\texttt{WAITAOF} returned
success for every durability acknowledgement requested after the device had
stopped accepting writes.} \texttt{dm-flakey}'s \texttt{drop\_writes} discards
each write bio and reports success, so the kernel sees no error,
\texttt{fsync(2)} succeeds, and Redis has no way to learn otherwise. The barrier
withholds dispatch on a \emph{failed or absent} durability acknowledgement; it
was handed a successful one and dispatched, behaving exactly as specified on an
input that was false. Both arms finished at ceiling.

We state the scope narrowly. \textbf{It is not evidence that the barrier fails
to withhold under write loss}: the condition the experiment exists to create ---
the barrier meeting an acknowledgement that fails or never arrives --- was never
reached, so no run exercised withholding. \textbf{Neither is it evidence that
the barrier works.} What it establishes is that the durability guarantee rests
on an unstated premise --- that a \texttt{WAITAOF} success means the record
reached the device --- and that under storage which discards writes
\emph{silently} the premise is false and the protocol cannot detect it. Storage
that returned \texttt{EIO} would let the acknowledgement fail and put the
barrier back on the input its guarantee is defined over; this result therefore
bounds the guarantee to storage that reports its own failures, and says nothing
about storage that does not. It is one host and one Redis version (7.2.5,
\texttt{appendfsync everysec}), and we claim it for no other.

\paragraph{The write-loss probe runs several layers above hardware, and does not
need to be closer} The \texttt{dm-flakey} target sits on a loop device backed by
a file on the WSL2 root filesystem, itself a virtual disk on the host. The
layering is irrelevant to the claim, because \texttt{dm-flakey} \emph{is} the
boundary under test: the experiment asks whether a \texttt{WAITAOF}
acknowledgement corresponds to bytes that have crossed a chosen point in the
storage path before that point stops accepting writes, and everything below the
target is on the far side of it in both arms equally. What the layering
\emph{would} threaten is a claim about absolute fsync latency on real hardware,
which we do not make from this probe --- the barrier's cost comes from
\cref{tab:deployment} --- or a claim that a real power cut cannot do worse,
which we explicitly do not make above.

\paragraph{No process-level fault can exercise the barrier's durability at all}
Stated separately because a reader is most likely to have assumed the other way.
\texttt{appendfsync everysec} defers the \texttt{fsync(2)} and not the
\texttt{write(2)}, so a \texttt{SIGKILL} of Redis loses nothing that
\texttt{appendonly yes} was not already keeping ($\ProcessKillUnackLost{}$, and
$n{=}\KillCanaryN{}$ canaries in the ablation cells). \textbf{Under a process
fault the barrier's contribution is prevention (\cref{sec:eval-prevention},
$p=\UnwantedP{}$, replicated at
[\ReplicationPreventedLow{},\,\ReplicationPreventedHigh{}] effects prevented)
and not durability}, and any reading of the barrier as protecting the record
from a crashed Redis process is wrong.

\paragraph{The barrier is evaluated under crash faults only} Every prevention
measurement here delivers an F3 crash-stop of the state store, a class in which
the arms are separated by a \emph{timing} difference and nothing else: AEP-full
has one extra Redis-dependent step, and the fault must land in a window B3 has
already left. That window can be narrowed and measured (\cref{sec:eval-prevention})
but not closed, because a fault that guarantees \texttt{WAITAOF} cannot return
also blocks the two Redis calls B3 makes after the same checkpoint. One
mechanism would avoid this --- disabling the fsync at the checkpoint while
leaving every other Redis call working, which separates the arms by construction
--- and \textbf{we designed that arm and did not run it}. It is a configuration
change that makes a guarantee unavailable rather than a fault, it would
instantiate a failure class our model does not contain, and the controlled fault
had already reduced the between-session spread to \ArmASpreadMax{} count per
class. The scope this fixes is explicit: the prevention guarantee is evaluated
where durability is \emph{raced for} and lost to a crash, not where it is simply
\emph{unavailable}. A reader who needs the latter should treat it as untested.

\paragraph{The mock provider's realism} The endpoint is a purpose-built service
with a configurable delay, timeout and error distribution, and a
transactionally written ground-truth ledger. It is realistic in the dimension
the experiment turns on --- its reconciliation capability, the independent
variable of \cref{tab:outcomes} --- and unrealistic in most others: no
authentication, no rate limits, no partial writes, no schema evolution, one
mutation shape. A real legacy endpoint may also \emph{misdeclare} its
capability, which the protocol cannot detect (\cref{sec:model-endpoint}); we
validate the capability's shape, never its truthfulness.

\paragraph{The dispatch guard assumes trusted same-process code} The guard
prevents accidental or ordinary control-flow bypass through the supported
dispatch API, and its registry detects copying, reuse and scope mismatch. It is
not a cryptographic capability: Python's underscore naming,
object-construction checks and closure placement are not security boundaries,
and code running in the same interpreter can reach module internals or bypass
the API altogether. The implementation therefore enforces a protocol invariant
only under a trusted-code assumption; a hostile-extension threat model would
require a separately enforced process or service boundary.

\paragraph{B4 is not Temporal} B4 shares exactly one mechanism with a production
durable-execution engine: durable history, replay, and at-least-once activity
re-execution. It does not model heartbeating, backoff timing, task queues,
versioning, or determinism checking. Its re-execution policy is the vendor's
documented default and is quoted as such~\cite{temporal-retry}, but the decision
is our code. No claim about Temporal's throughput, latency, maturity or
correctness as a product is made or supported anywhere in this paper. B4 and B4b
\emph{bracket} the configuration space rather than survey it: a deployment could
set Maximum Attempts to 3, add a heartbeat, or make the activity idempotent ---
the first two land between our two points on the same duplicate/loss trade, and
the third leaves the problem this paper is about.

\paragraph{The engine itself was run, and it separates what B4 licenses from
what it does not} We ran the production engine against the same non-cooperative
endpoint, in the two configurations B4 and B4b instantiate, at the one crash
point both the engine and our model can be cut at. With unlimited attempts it
duplicated at \BfiveDup{} over \BfiveExec{} executions and lost effects at
\BfiveLost{}; with Maximum Attempts set to one it lost at \BfivebLost{} over
\BfivebExec{} executions and duplicated at \BfivebDup{}. Neither configuration
declared ambiguity in any of \BfiveRuns{} non-void runs (\BfiveAmb{} and
\BfivebAmb{} declarations) --- the third corner of \cref{sec:trilemma} going
unreached by the engine itself and not merely by our model of it.

\emph{The direction transfers; the magnitudes do not.} Each configuration lands
in the corner \cref{tab:related} assigns it, with exactly zero events of the
other kind. But at that same crash point B4 duplicates at \BfourAtBarrier{} over
\BfourAtBarrierExec{} executions and B4b loses at \BfourbAtBarrier{} over
\BfourbAtBarrierExec{}, roughly six times the engine's rates. \textbf{B4's rates
are our model's, not the product's}, and no rate in this paper should be read as
an engine's.

Four things bound this and none is removed by having run the engine. The
comparison is \emph{incomplete}: the engine has no position answering to
\textsc{after-intent-before-barrier}, because its pre-dispatch record is written
by the server inside one RPC, so there is no worker-side instant between the
write and its acknowledgement --- a hole exactly where B4 and B4b have data at
all three capability classes. The rates are conditional on a
\BfiveStartToClose{}\,ms Start-To-Close timeout that B4 has no analogue for; a
different timeout is a different cell, and no timing number here is any engine's
recovery latency. The interval arithmetic is asymmetric --- B4's intervals rest
on three runs against the engine's thirty, and three clusters put resampled
means on multiples of one thirtieth --- so the separation should be read from
the point estimates, which differ by roughly $0.8$, and not from the intervals
failing to overlap. And the crash window at the dispatch point is sized by the
server and by loopback rather than by the protocol, so its width is not
like-for-like with B4's.

\paragraph{Declared ambiguity is not evaluated as an operational outcome} The
paper measures how often AEP declares ambiguity and shows that it never
substitutes a silent failure for it. It does \emph{not} show that declaring is
better in practice; that would require an operator study, and we have run none.
The argument that a declared incident beats an undiscovered one is a normative
claim about failure handling, not a result of this evaluation. Relatedly, the
artifact has no escalation \emph{mechanism} (\cref{tab:nonclaims}): reaching the
terminal state pauses the execution and alerts nobody.

\paragraph{Test counts are not protocol assurance} The suite's size is evidence
about the artifact's hygiene, not about the protocol's behaviour under faults.
The evidence for the latter is the fault-injection matrix.

\subsection{Internal validity}

\paragraph{The decomposition narrows our own novelty claim}
\Cref{sec:eval-detection} attributes the detection guarantee to the pre-dispatch
record and the transition table's missing re-entry edge, not to the barrier
machinery of contribution C2. The reading follows: our most novel mechanism
serves the claim with the \emph{narrowest} evidence --- one crash point, one
host, an effect size whose host-dependence we have measured rather than
suspected --- while the claim with the strongest evidence rests on write-ahead
logging plus a state-machine property that is a few lines of Lua. Narrowest is
not weakest: within that scope the effect replicated across
\ReplicationSessions{} further sessions at
[\ReplicationPreventedLow{},\,\ReplicationPreventedHigh{}] effects prevented,
the only session-clustered interval in this paper that excludes zero. The
composition being novel was never the argument (\cref{sec:related}); what is new
is the semantics it delivers, and the ablation says which half delivers which
--- including that the expensive half is optional.

The cell is no longer one collection under one injector.
\Cref{sec:eval-prevention} leads with a \emph{controlled} fault --- the
container frozen before it is killed, landing in \ControlledLanding{}\,ms
against \UncontrolledLanding{}\,ms --- over \ArmASessions{} sessions and all
\ArmAClasses{} capability classes. AEP-full's unwanted-applied rate is
\ArmAAepRate{} against B3's \ArmABthreeRate{}, with a between-session spread of
\ArmASpreadMax{} count per class: a measured quantity with a stated spread
rather than one draw from a distribution we could only name. \textbf{It is
bounded by measurement rather than by prediction, and that cost us a
pre-registration.} We registered [\ArmABoundLow{},\,\ArmABoundHigh{}] before
collecting, derived from the injector's landing latency as a fraction of the
\texttt{WAITAOF} window, and the observed \ArmAAepRate{} fell \emph{below} it.
Both of the model's premises fail on measurement: \ModelAcksPastLanding{} of
\ModelAckN{} uncontrolled runs acknowledged after that landing had supposedly
passed and the same injector call runs \ModelInSituRatio{}$\times$ the bench
figure in situ, and the wait is not uniform --- \ModelEmpiricalWindow{} of waits
fall under \ControlledLanding{}\,ms where uniformity assigns
\ModelUniformWindow{}. Take the measured rate and its spread, not the bound.

\textbf{The residual is real and irreducible under this fault class.} AEP-full
still applied \ArmAAepApplied{} effects in \ArmAAepN{} executions, and no
crash-class fault can drive that to zero, for the reason given under construct
validity. The uncontrolled cell agrees in direction and its magnitude still
moves: \ReplicationSessions{} collections ranged
\ReplicationAepMin{}--\ReplicationAepMax{} applied effects out of
\AepKillRuns{}. Attributing that movement to kill latency did not succeed at
this design's precision --- \KillLatencyPerSession{}\,ms per session, mean
\KillLatencyMean{}\,ms, 95\% interval
[\KillLatencyLow{},\,\KillLatencyHigh{}], a half-width
\KillLatencyPrecisionRatio{} times that mean, containing zero because the
sessions disagree. \UnwantedPrevented{} is one draw from that distribution;
\ArmAAepRate{} is not.

\paragraph{The detection finding has no referent outside this artifact} The
decomposition above is established by an ablation --- AEP-full against itself
with the barrier removed --- and an ablation is internal by construction. That
is the right instrument for \emph{attributing} an effect to a mechanism, but it
cannot corroborate the resulting claim from outside. Nothing external to this
artifact establishes that a pre-dispatch record without an fsync barrier
suffices for detection: the proposition is supported by two systems we wrote,
measured by a harness we wrote, against a provider we wrote. The oracle's
independence (below) addresses whether we counted correctly, not whether the
finding generalises beyond our implementation of the pattern. That is the
structural cost of the instrument, and it applies to the paper's most
load-bearing empirical claim.

\paragraph{The matrix found defects in our own measurement, including one that
would have inflated counts} Three are on the record. A resume defect let a
re-attempted run merge stale event shards into its own log, double-counting
executions. A re-executing supervisor abandoned the lease instead of waiting for
it, and execution-state seeding was not idempotent --- both making a
\emph{baseline} look better than it is, and so biasing against our own
hypothesis. And running the project's own destructive integration suite beside a
live collection replaced the matrix's Redis underneath it, so every B4/B4b run
from that window was discarded and re-collected. All three are fixed, the
affected runs were re-collected, and none of the discarded runs contributes a
number here; the runner now refuses to start beside a host-level fault injector
and records the server's \texttt{run\_id} with every run. The supplementary
material gives each incident in full. We record them because the general point
--- that a fault-injection harness is a stronger verification tool than a unit
suite --- cuts both ways, and there may be defects it has not found.

\paragraph{Ambiguity is counted from the oracle, not from self-report} A system
that declared ambiguity liberally could otherwise score well by saying ``I do
not know'' about everything. Duplicates and lost effects are counted by
comparing the provider's independently written ledger against the system's
records; the analysis is prevented by a parsed source gate from reading live
protocol state.

\paragraph{The oracle is independent, not infallible} The provider ledger is
written outside the systems under test, and every completed run is reconciled
against it, but the provider and ledger share one implementation and one host.
One voided run demonstrates that this oracle can fail: its event log and ledger
disagreed while sibling and repeated runs agreed. The current repository retains
the aggregate record of that exclusion, and the voided raw directory must be
included in the external archive before submission. A second provider
implementation or an external observer would reduce this common-mode doubt;
neither is present.

\paragraph{One author implemented every system under comparison} B0--B2 are
small enough to inspect, each carries a machine-readable descriptor that tests
check against its implementation, and B4's re-execution policy is defended
against the vendor's own documentation in the artifact. That is mitigation, not
elimination: a baseline written by the people proposing the alternative is a
structural conflict, and the artifact is released partly so the baselines can be
re-implemented by someone else.

\subsection{External validity}

\paragraph{Single node, single trust domain} One Redis instance, no replica, no
consensus. \Cref{tab:nonclaims} lists what this forecloses. One residual
interacts with the results: an AOF rewind can un-fence a lease, since lock keys
are deliberately not barriered, and the fix is HA rather than anything local.

\paragraph{Concurrency coverage is narrow} The harness exercises crash and
takeover with two workers, so stale-owner and successor paths are covered, but
it does not sustain live contention among many workers for the same execution or
endpoint resource. The evaluation therefore supports fencing and recovery under
the measured takeover schedule, not a scalability or fairness claim for
high-contention deployments. Concurrent-claim rejection is exercised by the test
suite; the fault-injection matrix supplies no contention-regime evidence for P1.

\paragraph{Platform, and a development/measurement split} All measurements were
collected on Ubuntu inside WSL2 on Windows~11, with Docker Desktop port
forwarding in the path, on mains power with sleep and hibernation disabled.
Development and part of the test suite run on the Windows host; \emph{every}
number in \cref{sec:evaluation} comes from the Linux tree against real Redis.
The locally runnable Windows suite and the artifact's Linux suite are checked
separately, while the Linux evidence is the one the measurements rest on.
Absolute latency and throughput are platform-bound; only comparisons
\emph{between} systems sharing this platform are meaningful, which is why
\cref{tab:latency} is presented as increments over a common floor.

\paragraph{The motivating traces come from our own harness} \Cref{sec:motivating}
reproduces failures in baselines we implemented, not in deployed systems. It
establishes that the failure modes are reachable under the stated fault model;
it is not evidence about their frequency in production deployments, and we have
none. A field study of how often such systems meet non-idempotent endpoints
without idempotency keys would strengthen the motivation considerably.

\paragraph{Coverage of the matrix is partial, and the gaps are named} The full
plan is 1\,068 runs; we collected \RunsCollected{}, comprising
\ExecutionsCollected{} executions over \CellsCollected{} cells, chosen to answer
the research questions rather than to fill the grid.
% source: experiments/results/matrix/MANIFEST.md, which lists the run count
% for every cell, and analysis/coverage.json.
The supplementary material enumerates six named gaps --- three since closed and
recorded there with what closing them changed, three still open: the $30\%$
crash-probability regime, the alternative read-back keying, and one run that did
not complete. Two qualifications from that enumeration bear on the main results
and are stated here rather than only there: every collection in this paper comes
from \emph{one host}, so the prevention effect's dependence on host timing is
measured \emph{on} that host rather than established \emph{across} hosts; and
completing a previously partial cell changed its value materially, so a partial
cell should be read as unrepresentative rather than merely imprecise. A rate we
did not measure is absent from the tables rather than interpolated, and the
manifest shipped with the results states the run count for every cell.

\paragraph{One endpoint for the cost measurement} RQ3's cells were collected
against a single endpoint. Overhead should not depend on an endpoint's
reconciliation capability in a crash-free run, since no read-back occurs, but
that is an argument and not a measurement.

\subsection{Conclusion validity}

\paragraph{Every timing number rests on three runs, and one does not survive its
own interval} Each cell behind \cref{tab:deployment} is three crash-free runs of
ten executions, and the four medians there are what the barrier's cost is
computed from. A cluster bootstrap over runs (\cref{sec:eval-deployment}) changes
what can be claimed. Under \texttt{everysec} the barrier's cost is
\BarrierCost{}\,ms with a 95\% interval of
[\BarrierCostLow{}, \BarrierCostHigh{}] --- positive, but pinned only to within
about a factor of four. Under \texttt{always} the point estimate is
\BarrierCostAlways{}\,ms and the interval [\BarrierCostAlwaysLow{},
\BarrierCostAlwaysHigh{}] spans zero: at three runs per arm we cannot show the
barrier costs anything under that policy. An earlier draft quoted a factor
between the two point estimates; it is not supported by these data and has been
removed. The cause is the sample, not the metric: three clusters admit only ten
distinct bootstrap multisets, and the \texttt{always} cell has a heavy upper tail
(p95 \AepAlwaysPninetyfive{}\,ms against a \AepAlwaysMedian{}\,ms median) that
three clusters cannot resolve. Closing this needs more crash-free \emph{runs},
not more executions per run. A deployment that needs a latency budget should
measure its own.

\paragraph{A failure to reject is not a null result here} Two comparisons in
this paper fail to exclude zero for reasons of precision rather than of effect.
The kill-latency attribution above has a half-width wider than the mean it
brackets. So does the capability-class sweep reported in the supplementary
material, where the applied-effect column moved --- contradicting our own
registered prediction --- while the session-clustered interval on the difference
still contains zero because the sessions disagree with one another. In both
cases what we report is a failure to reject at a precision inadequate to the
question, and neither should be read as evidence that the effect is absent.

\paragraph{Timing hygiene, and what it cost us} The host's modern-standby
setting suspended it mid-run during earlier collection, silently inflating
wall-clock durations. We now require both a pre-run declaration that the host
cannot suspend and a per-run wall-versus-monotonic check. The gate is strict
enough that \emph{no} run collected before it existed contributes a timing
number, and it costs us coverage: recovery-latency distributions are not
reported (\cref{sec:eval-rq4}) because the gated crashed runs are too few, and
we prefer the gap to ungated numbers.

\subsection{Where the protocol should not be used}

If the endpoint accepts an idempotency key, or can be enrolled in a
transaction, or the effect can be made idempotent, then a durable-execution
engine gives a strictly stronger guarantee at lower cost and AEP is the wrong
tool. Its cost is real: two fsync barriers, which under \texttt{appendfsync
everysec} dominate the protocol's latency by a factor of roughly
\BarrierToProtocolRatio{} (\cref{tab:latency}) --- a ratio of medians, and one
whose denominator is pinned only to
[\ProtocolMinusBarrierLow{}, \ProtocolMinusBarrierHigh{}]\,ms at three runs per
arm, so the factor is an estimate of an order of magnitude rather than of a
figure. AEP is for the case where the precondition those systems
require cannot be obtained, and its value is a declared residual rather than a
better guarantee.


% IEEE Editorial Style Manual for Authors: the acknowledgment goes after
% the final text and any appendixes and BEFORE the references; its heading
% is never enumerated; and the spelling is always the singular
% "Acknowledgment", with no "e". Hence \section*, here, spelled this way.
\section*{Acknowledgment}
\aidisclosure

\bibliographystyle{IEEEtran}
\bibliography{refs}

\end{document}
